\documentclass{amsart}
\usepackage{amsmath, amsthm, amsfonts, amssymb, color,enumerate}
\usepackage[colorlinks, citecolor=blue,pagebackref,hypertexnames=false]{hyperref}


\begin{document}

\newtheorem{theorem}{Theorem}[section]
\newtheorem{proposition}[theorem]{Proposition}
\newtheorem{coro}[theorem]{Corollary}
\newtheorem{lemma}[theorem]{Lemma}
\newtheorem{definition}[theorem]{Definition}
\newtheorem{nota}[theorem]{Notation}
\newtheorem{assum}{Assumption}[section]
\newtheorem{example}[theorem]{Example}
\newtheorem{remark}[theorem]{Remark}


\title[Schatten class properties of Dunkl--Riesz transform commutators]{Weyl--Chamber geometry, Poincar\'e inequality meets Schatten class properties of Dunkl--Riesz transform commutators}



\author{Zhijie Fan}
\address{Zhijie Fan, School of Mathematical Sciences, South China Normal University, Guangzhou 510631, China}
\email{fanzhj3@mail2.sysu.edu.cn}


\author{Ji Li}
\address{Ji Li, Department of Mathematics, Macquarie University,NSW2109, Australia}
\email{ji.li@mq.edu.au}



\author{Dmitriy Zanin}
\address{School of Mathematics and Statistics, HNP-LAMA, Central South University, Changsha 410075, China}
\email{d.zanin@csu.edu.cn}






\date{}

 \subjclass[2010]{47B10, 42B20, 43A85}
\keywords{Schatten class, Riesz transform commutators, Dunkl operator, Besov space, Nearly weakly orthogonal sequence}




\begin{abstract}

\end{abstract}

\maketitle


\tableofcontents


\section{Introduction}

\begin{theorem}\label{main theorem} Let $R$ be a normalised root system in $\mathbb{R}^N,$ $N\geq 2,$ and let $G$ be the corresponding finite reflection group. Let $k:R\to\mathbb{R}_+$ be $G$-invariant multiplicity function. Let $R_{{\rm Dunkl},j}$ be the $j$-th Dunkl-Riesz transform corresponding to $R$ and $k.$ Let $w$ be the Dunkl weight corresponding to $R$ and $k.$ Let $f\in L_1^{{\rm loc}}(\mathbb{R}^N).$
\begin{enumerate}[{\rm (i)}]
\item\label{mta} If $f\in\dot{W}^{1,N}(\mathbb{R}^N),$ then $[R_{{\rm Dunkl},j},M_f]\in\mathcal{L}_{N,\infty}(L_2(\mathbb{R}^N,w(x)dx)).$ Furthermore,
$$\|[R_{{\rm Dunkl},j},M_f]\|_{\mathcal{L}_{N,\infty}(L_2(\mathbb{R}^N,w(x)dx))}\leq c_{R,k}\|f\|_{\dot{W}^{1,N}(\mathbb{R}^N)};$$
\item\label{mtb} If $[R_{{\rm Dunkl},j},M_f]\in\mathcal{L}_{N,\infty}(L_2(\mathbb{R}^N,w(x)dx)),$ then $f\in\dot{W}^{1,N}(\mathbb{R}^N).$ Furthermore,
$$\|[R_{{\rm Dunkl},j},M_f]\|_{\mathcal{L}_{N,\infty}(L_2(\mathbb{R}^N,w(x)dx))}\geq c_{R,k}^{-1}\|f\|_{\dot{W}^{1,N}(\mathbb{R}^N)};$$
\end{enumerate}
\end{theorem}

\section{Preliminaries}
\setcounter{equation}{0}
\subsection{Harmonic analysis in the Dunkl setting}
Let $\mathbb{R}^N$ be the Euclidean space equipped with the usual scalar product $\langle \cdot,\,\cdot\rangle $ and the induced norm $\|\cdot\|$. For a nonzero vector $v\in \mathbb{R}^N$, the reflection $\sigma_v$ with respect to the hyperplane orthogonal to a normalised vector $v$ is defined by
$$\sigma_v x:= x-2\langle x,v\rangle v.$$
A finite set $R\subset \mathbb{R}^N\backslash \{0\}$ is called a root system if $\sigma_v(R)=R$ for every $v\in R.$ A root system $R$ is called a normalised reduced root system if  $\|v\|^2=1$ for every $v\in R.$
The finite group $G$ generated by the reflections $\{\sigma_v\}_{v\in R}$ is called the Weyl group (reflection group) of the root system $R$.
Let $\kappa: R \to [0,\infty)$ be a multiplicity function, which is a $G$-invariant function  fixed throughout the paper. 


Define $d(x,y)=\min_{g\in G}|x-gy|.$ Recall from \cite[Chapter VII, proof of Theorem 2.12]{MR514561} that $d(x,y)=|x-y|$ iff $x$ and $y$ belong to the same chamber.

Define the Dunkl measure $dw$ on $\mathbb{R}^N$ by
$$
dw(x) :=\prod_{v \in R} |\langle x,v\rangle|^{\kappa(v)} dx,
$$
where $dx$ stands for the Lebesgue measure on $\mathbb{R}^N$. Let $B(x,r):=\{y\in \mathbb{R}^N:\|x-y\|< r\}$ be the Euclidean ball centred at $x$ with radius $r>0$. Let $\mathbf{N}:=N+\sum_{v\in R}k(v)$. Note that
$$w(B(tx,tr))=t^{\mathbf{N}}w(B(x,r)),\ {\rm for}\ x\in\mathbb{R}^N,\ t,r>0.$$
Furthermore,
\begin{align}\label{twosideinequ}
  w(B(x,r))\sim r^N \prod_{v \in R}(|\langle x, v \rangle |+r)^{\kappa(v)},
\end{align}
where the implicit constant only depends on $N,R$ and $\kappa.$  It follows from \eqref{twosideinequ} that there exists a constant $C\geq 1$ such that for every $x\in \mathbb{R}^N$,
\begin{align}\label{doublinginequality}
C^{-1}\left(\frac{r_2}{r_1}\right)^N\leq \frac{w(B(x,r_2))}{w(B(x,r_1))}\leq C\left(\frac{r_2}{r_1}\right)^{\mathbf{N}},\ {\rm for}\ 0<r_1<r_2,
\end{align}
so the numbers $\mathbf{N}$ and $N$ are called the upper and lower homogeneous dimensions, respectively.

The Dunkl operators $T_\xi$ are defined by
$$T_\xi f(x):=\partial_\xi f(x)+\sum_{v\in R}\frac{k(v)}{2}\langle v,\xi\rangle\frac{f(x)-f(\sigma_v(x))}{\langle v,x\rangle},$$
where $\partial_\xi f$ denotes the directional derivative of $f$ in the direction $\xi$. Furthermore, the Dunkl Laplacian is defined by
$$\Delta_{{\rm Dunkl}}=-\sum_{j=1}^N T_{e_j}^2.$$

In this article, we consider the Dunkl-Riesz transforms defined by $R_{{\rm Dunkl},j}=T_{e_j}\Delta_{{\rm Dunkl}}^{-\frac12},$ $j=1,2,\cdots,N.$

\subsection{Weyl chamber as simplicial cone}

\begin{definition} A connected component of the set
$$\left\{x\in\mathbb{R}^N:\langle x,v\rangle \neq 0\ \ { for} \ { all} \ v\in R\right\}$$
is called an open Weyl chamber. Its closure is called a Weyl chamber.
\end{definition}

The following assertion can be found in e.g. Theorem 2 on p.166 in \cite{Bourbaki}.

\begin{theorem}\label{phi theorem} Let $\Gamma$ be a chamber. There exists a linear isomorphism $\Phi:\mathbb{R}^N\to\mathbb{R}^N$ such that $\Phi(\Gamma)=R_+^{N_1}\times\mathbb{R}^{N_2}$ (with $N_1+N_2=N$).
\end{theorem}




\subsection{System of dyadic cubes}

\begin{nota}\label{standardsystem}
Let
$$\mathbb{D}=\bigcup_{k\in\mathbb{Z}}\mathbb{D}_k$$
be the standard system of dyadic cubes over $\mathbb{R}^N.$ Here each dyadic cube in $\mathbb{D}_k$ is of the shape
$$\frac{m}{2^k}+[0,2^{-k})^N,\quad m\in\mathbb{Z}^N.$$

For every dyadic cube $Q\in\mathbb{D}$, we denote by $c_Q$ and $\ell(Q)$ its center and side-length, respectively.
\end{nota}


\subsection{Nearly weakly  orthonormal sequences}
Rochberg and Semmes \cite{RS} introduced the notion of \emph{nearly weakly  orthonormal (NWO)} sequences of functions. This notion has since proved highly effective in characterizing Schatten class membership of  singular integral commutators in a variety of contexts (see \cite{HKarxiv,MR4756021,MR4552581,MR4873796,LLW,RS,Zhangwei1,Zhangwei2}). The properties of NWO sequences were originally investigated in the Euclidean setting (equipped with Lebesgue measure), but the extension to spaces of homogeneous type is straightforward once a system of dyadic cubes is available in this general framework (see \cite[Section 13]{HKarxiv} for a complete proof).

{\color{red} Definition \ref{NWO def} and Lemma \ref{NWO lemma} as stated below are my understanding of Hytonen's Corollary 8.2 and Lemma 8.7. Please check whether this understanding is the correct one!}

\begin{definition}\label{NWO def} Let ${\bf w}\in A_2(\mathbb{R}^N).$ Let $c>0$ and $2<r\leq\infty.$ A sequence of functions $\{e_Q\}_{Q\in\mathbb{D}}$ is called $(r,c)$-NWO if
\begin{enumerate}[{\rm (i)}]
\item ${\rm supp}(e_Q)\subseteq cQ$ for every $Q\in\mathbb{D};$ 
\item $\|e_Q\|_{L_r(\mathbb{R}^N,{\bf w}(x)dx)}\leq {\bf w}(Q)^{\frac1r-\frac12}$ for every $Q\in\mathbb{D};$
\end{enumerate}
\end{definition}

\begin{lemma}\label{NWO lemma} Let ${\bf w}\in A_2(\mathbb{R}^N).$ Let $0<p<\infty.$ Suppose that $A$ is an integral operator on $L_2(\mathbb{R}^N,{\bf w}(x)dx)$ with integral kernel 
$$(x,y)\to \sum_{Q\in\mathbb{D}}\lambda_Q e_Q(x)f_Q(y),\quad x,y\in\mathbb{R}^N,$$
where $\{e_Q\}_{Q\in\mathbb{D}}$ and $\{f_Q\}_{Q\in\mathbb{D}}$ are $(r_1,c_1)$-NWO and $(r_2,c_2)$-NWO sequences, respectively. We have
$$\|A\|_{\mathcal{L}_{p,\infty}(L_2(\mathbb{R}^N,{\bf w}(x)dx))}\leq C_{w,r_1,c_1,r_2,c_2,p} \|\{\lambda_Q\}_{Q\in\mathbb{D}}\|_{p,\infty},$$
$${\rm dist}_{\mathcal{L}_{p,\infty}(L_2(\Gamma,{\bf w}(x)dx))}\Big(A,(\mathcal{L}_{p,\infty})_0(L_2(\mathbb{R}^N,{\bf w}(x)dx))\Big)\leq $$
$$\leq C_{{\bf w},r_1,c_1,r_2,c_2,p}{\rm dist}_{l_{p,\infty}(\mathbb{D})}\Big(\{\lambda_Q\}_{Q\in\mathbb{D}},(l_{p,\infty})_0(\mathbb{D})\Big).$$
\end{lemma}
\begin{proof} The assertion follows from Corollary 8.2 and Lemma 8.7 in \cite{HKarxiv}.
\end{proof}

\begin{lemma}\label{NWO lower estimate} Let $w\in A_2(\mathbb{R}^N).$ Let $1<p<\infty.$ Suppose that $A$ is an integral operator on $L_2(\mathbb{R}^N,{\bf w}(x)dx)$ with integral kernel 
$$(x,y)\to \sum_{Q\in\mathbb{D}}\lambda_Q e_Q(x)f_Q(y),\quad x,y\in\mathbb{R}^N,$$
where $\{e_Q\}_{Q\in\mathbb{D}}$ and $\{f_Q\}_{Q\in\mathbb{D}}$ are $(r_1,c_1)$-NWO and $(r_2,c_2)$-NWO sequences, respectively. We have
$$\Big\|\Big\{\langle A e_Q,f_Q \rangle_{L_2(\mathbb{R}^N,{\bf w}(x)dx)}\Big\}_{Q\in\mathbb{D}}\Big\|_{l_{p,q}(\mathbb{D})}\leq c_{{\bf w},r_1,c_1,r_2,c_2,p,q}\|A\|_{\mathcal{L}_{p,q}(L_2(\mathbb{R}^N,{\bf w}(x)dx))}.$$
\end{lemma}
\begin{proof} The assertion follows from Corollary 8.2 and Lemma 8.7 in \cite{HKarxiv}.
\end{proof}




\section{Chamber is a doubling metric measure space}

\begin{theorem}\label{main:dunkldoubling} Let $\Gamma$ be a chamber. For all $x\in\Gamma$ and for all $r>0,$ we have
$$w(B(x,r)\cap\Gamma)\geq c_{\kappa}w(B(x,r)).$$
In particular, $(\Gamma,|\cdot|,w(x)dx)$ is a doubling metric measure space.
\end{theorem}

\begin{nota} Let $\theta:2^{N_1}\to\mathbb{R}_+.$ Set
$$w_{\theta}(x)=\prod_{A\subset\{1,\cdots,N_1\}}(\sum_{k\in A}x_k)^{\theta_A},\quad x\in\mathbb{R}_+^{N_1}.$$
\end{nota}

\begin{lemma}\label{s2 upper estimate} For every $\theta:2^{N_1}\to\mathbb{R}_+$ and for every $x\in\mathbb{R}^{N_1}_+\times\mathbb{R}^{N_2},$ we have
$$\int_{B(x,r)}w_{\theta}(y_1,\cdots,y_{N_1})dy\leq c_{N,\theta}^{(1)}r^Nw_{\theta}((x_1,\cdots,x_{N_1})+r\mathbf{1}_{N_1}).$$
\end{lemma}
\begin{proof} Obviously,
$$B(x,r)\subset \Big(\times_{k=1}^{N_1}[(x_k-r)_+,x_k+r]\Big)\times B((x_{N_1+1},\cdots,x_N),r).$$
Thus,
$$\int_{B(x,r)}w_{\theta}(y_1,\cdots,y_{N_1})dy\leq$$
$$\leq\int_{\big(\times_{k=1}^{N_1}[(x_k-r)_+,x_k+r]\big)\times B((x_{N_1+1},\cdots,x_N),r)}w_{\theta}(y_1,\cdots,y_{N_1})dy=$$
$$=\int_{\times_{k=1}^{N_1}[(x_k-r)_+,x_k+r]}w_{\theta}(u)du\times {\rm Vol}(\mathbb{B}_{N_2})r^{N_2}\leq$$
$$\leq\int_{\times_{k=1}^{N_1}[0,x_k+r]}w_{\theta}((x_1,\cdots,x_{N_1})+r\mathbf{1}_{N_1})du\times {\rm Vol}(\mathbb{B}_{N_2})r^{N_2}=$$
$$=w_{\theta}((x_1,\cdots,x_{N_1})+r\mathbf{1}_{N_1})\times\prod_{k=1}^{N_1}(x_k+r-(x_k-r)_+)du\times {\rm Vol}(\mathbb{B}_{N_2})r^{N_2}\leq $$
$$\leq w_{\theta}((x_1,\cdots,x_{N_1})+r\mathbf{1}_{N_1})\times(2r)^{N_1}\times {\rm Vol}(\mathbb{B}_{N_2})r^{N_2}.$$
\end{proof}

\begin{lemma}\label{s2 lower estimate} For every $\theta:2^{N_1}\to\mathbb{R}_+$ and for every $x\in\mathbb{R}^{N_1}_+\times\mathbb{R}^{N_2},$ we have
$$\int_{B(x,r)\cap(\mathbb{R}^{N_1}_+\times\mathbb{R}^{N_2})}w_{\theta}(y_1,\cdots,y_{N_1})dy\geq c_{N,\theta}^{(2)}r^Nw_{\theta}((x_1,\cdots,x_{N_1})+r\mathbf{1}_{N_1}).$$
\end{lemma}
\begin{proof} Obviously,
$$x+[\frac{r}{2N},\frac{r}{N}]^N\subset B(x,r)\cap (\mathbb{R}^{N_1}_+\times\mathbb{R}^{N_2}).$$
It follows that
$$\int_{B(x,r)\cap(\mathbb{R}^{N_1}_+\times\mathbb{R}^{N_2})}w_{\theta}(y_1,\cdots,y_{N_1})dy\geq $$
$$\geq\int_{x+[\frac{r}{2N},\frac{r}{N}]^N}w_{\theta}(y_1,\cdots,y_{N_1})dy=$$
$$=\int_{(x_1,\cdots,x_{N_1})+[\frac{r}{2N},\frac{r}{N}]^{N_1}}w_{\theta}(u)du\times (\frac{r}{2N})^{N_2}\geq $$
$$\geq \int_{(x_1,\cdots,x_{N_1})+[\frac{r}{2N},\frac{r}{N}]^{N_1}}w_{\theta}((x_1,\cdots,x_{N_1})+\frac{r}{2N}\mathbf{1}_{N_1})du\times (\frac{r}{2N})^{N_2}=$$
$$=w_{\theta}((x_1,\cdots,x_{N_1})+\frac{r}{2N}\mathbf{1}_{N_1})\times(\frac{r}{2N})^N\geq $$
$$\geq w_{\theta}((x_1,\cdots,x_{N_1})+r\mathbf{1}_{N_1})\times(\frac{r}{2N})^N\times (2N)^{-\sum_{A\subset\{1,\cdots,N_1\}}\theta_A}.$$
\end{proof}


\begin{lemma}\label{s2 trivial lemma} There exists $\theta:2^{N_1}\to\mathbb{R}_+$ such that
$$c_{\kappa}^{-1}w_{\theta}(x_1,\cdots,x_{N_1})\leq w(\Phi^{-1}(x))\leq c_{\kappa}w_{\theta}(x_1,\cdots,x_{N_1}),\quad x\in\mathbb{R}_+^{N_1}\times\mathbb{R}^{N_2}.$$
\end{lemma}
\begin{proof} For every $x\in\mathbb{R}^{N_1}_+\times\mathbb{R}^{N_2}$ and for every $v\in R,$ we have
$$\langle \Phi^{-1}(x),v\rangle=\langle x,(\Phi^{-1})^{\ast}(v)\rangle.$$
The latter expression has the same sign (positive or negative) on the whole cone $\mathbb{R}^{N_1}_+\times\mathbb{R}^{N_2}.$ Hence,
$$\big((\Phi^{-1})^{\ast}(v)\big)_k=
\begin{cases}
\mbox{something positive},& 1\leq k\leq N_1\\
0,& k>N_1
\end{cases}
$$
Let $A_v$ be the set of those $1\leq k\leq N_1$ with $\big((\Phi^{-1})^{\ast}(v)\big)_l$ being strictly positive. We, obviously, have
$$c_{\Phi,v}^{-1}\sum_{k\in A_v}x_k\leq |\langle \Phi^{-1}(x),v\rangle|\leq c_{\Phi,v}\sum_{k\in A_v}x_k,\quad x\in \mathbb{R}^{N_1}_+\times\mathbb{R}^{N_2}.$$
Thus,
$$\prod_{v\in R}c_{\Phi,v}^{-\kappa(v)}(\sum_{k\in A_v}x_k)^{\kappa(v)}\leq w(x)\leq \prod_{v\in R}c_{\Phi,v}^{\kappa(v)}(\sum_{k\in A_v}x_k)^{\kappa(v)},\quad x\in\mathbb{R}_+^{N_1}\times\mathbb{R}^{N_2}.$$
This completes the proof.
\end{proof}

\begin{proof}[Proof of Theorem \ref{main:dunkldoubling}] The first assertion follows by combining Lemmas \ref{s2 upper estimate}, \ref{s2 lower estimate} and \ref{s2 trivial lemma}. The second assertion is an obvious consequence of the first one.
\end{proof}


\section{Rochberg-Semmes based estimate}

\begin{theorem}\label{RS-type theorem} Let $N\geq 2,$ $1\leq N_1\leq N$ and $N_2=N-N_1.$ Let $\Omega=\mathbb{R}^{N_1}_+\times\mathbb{R}^{N_2}.$ Let ${\bf w}\in A_2(\Omega)$ be such that
$$\sup_{x\in Q}{\bf w}(x)\leq c_w\frac{{\bf w}(Q)}{m(Q)},\quad Q\in\mathbb{D},\quad Q\subset\Omega.$$
Suppose that $K:\Omega\times\Omega\to\mathbb{C}$
satisfies
$$\bigl|\partial_x^{\alpha}\partial_y^{\beta}K\bigr|(x,y)\leq c_{K,\alpha,\beta}\frac{|x-y|^{1-|\alpha|_1-|\beta|_1}}{{\bf w}(B(x,|x-y|))},\quad x,y\in\Omega,\quad \alpha,\beta\in\mathbb{Z}^N_+.$$
Let $V$ be an integral operator on $L_2(\Omega,{\bf w}(x)dx)$ with integral kernel $K.$ Suppose $N>2s.$ For every $f\in C^{\infty}_c(\mathbb{R}^N),$ we have
$$\|M_fV\|_{\mathcal{L}_{N,\infty}(L_2(\Omega,{\bf w}(x)dx))}\leq c_{K,{\bf w}}\| f\|_{L_N(\Omega)}.$$
\end{theorem}

\begin{lemma} There exists a smooth function $\psi$ supported on $[-1,1]^N\backslash(-\frac14,\frac14)^N$ such that
$$1=\sum_{Q\in\mathbb{D}}\chi_Q(x)\chi_{3Q}(y)\psi(\frac{x-y}{\ell(Q)}).$$
\end{lemma}
\begin{proof} Choose $\rho\in C_c^{\infty}(\mathbb{R}^N)$ such that $\rho=1$ on $(-\frac12,\frac12)^N$ and $\rho=0$ outside $(-1,1)^N.$ Set $\psi(z)=\rho(z)-\rho(2z).$ Clearly,
$${\rm supp}(\psi)\subset
\left\{z\in\mathbb{R}^N:\quad 
\frac14\leq |z|_{\infty}\leq 1\right\}.$$
For every $z\neq0,$ we have
\begin{equation}\label{eq:rankone-psi-partition}
\sum_{j\in\mathbb{Z}}\psi(2^jz)=1.
\end{equation}

For every $x,$ there exists a unique $Q\in\mathbb{D}_j$ containing $x.$ For this $Q,$ we have $\ell(Q)=2^{-j}.$ We have
$$1=\sum_{j\in\mathbb{Z}}\psi(2^j(x-y))=\sum_{j\in\mathbb{Z}}(\sum_{Q\in\mathbb{D}_j}\chi_Q(x))\psi(2^j(x-y))=$$
$$=\sum_{j\in\mathbb{Z}}\sum_{Q\in\mathbb{D}_j}\chi_Q(x)\psi(\frac{x-y}{\ell(Q)})=\sum_{Q\in\mathbb{D}}\chi_Q(x)\psi(\frac{x-y}{\ell(Q)}).$$

If $x\in Q$ and $\psi(\frac{x-y}{\ell_Q})\neq 0,$ then
$$x\in Q,\quad |x-y|_{\infty}\leq\ell(Q).$$
We obviously have $|x-c_Q|_{\infty}\leq\frac12\ell(Q)$ and $|x-y|_{\infty}\leq\ell(Q).$ Thus, $|y-c_Q|_{\infty}\leq\frac32\ell(Q).$ In other words, $y\in 3Q.$ Thus,
$$\chi_Q(x)\psi(\frac{x-y}{\ell(Q)})=\chi_Q(x)\chi_{3Q}(y)\psi(\frac{x-y}{\ell(Q)}).$$
This completes the proof.
\end{proof}

\begin{lemma}\label{lem:Weyl-off-diagonal-kernel-NWO} Let $r>2.$ In the setting of Theorem \ref{RS-type theorem}, we have
$$f(x)K(x,y)=\sum_{k,l\in\mathbb{Z}^N}(1+|k|^2+|l|^2)^{-N-1}\sum_{\substack{Q\in\mathbb{D}\\ Q\subset\Omega}}e_{k,l,Q}(x)f_{k,l,Q}(y),$$
where
\begin{enumerate}[{\rm (i)}]
\item $e_{k,l,Q}$ is supported in $Q$ and $f_{k,l,Q}$ is supported in $3Q$ for every $Q\in\mathbb{D}$ with $Q\subset\Omega;$
\item for every $k,l\in\mathbb{Z}^N$ and for every $Q\in\mathbb{D}$ with $Q\subset\Omega,$ we have
$$\|e_{k,l,Q}\|_{L_r(\mathbb{R}^N,{\bf w}(x)dx)}=\|f\|_{L_r(Q,{\bf w}(x)dx)};$$
\item for every $k,l\in\mathbb{Z}^N$ and for every $Q\in\mathbb{D}$ with $Q\subset\Omega,$ we have
$$\|f_{k,l,Q}\|_{L_{\infty}(\mathbb{R}^N)}\leq \frac{c_{K,\theta}\cdot\ell(Q)}{{\bf w}(Q)};$$
\end{enumerate}
\end{lemma}
\begin{proof} For every $x,y\in \Omega,$ we have
$$K(x,y)=\sum_{Q\in\mathbb{D}}\chi_Q(x)\chi_{3Q}(y)K(x,y)\psi(\frac{x-y}{\ell(Q)})=$$
$$=\sum_{\substack{Q\in\mathbb{D}\\ Q\subset\Omega}}\chi_Q(x)\chi_{3Q}(y)K(x,y)\psi(\frac{x-y}{\ell(Q)})=$$

For every $Q\in\mathbb{D}$ with $3Q\subset\Omega,$ define the map
$$K_Q:[-\frac12,\frac12]^N\times[-\frac32,\frac32]^N\to\mathbb{C}$$
by setting
$$K_Q(a,b)=K(c_Q+a\ell(Q),c_Q+b\ell(Q))\psi(a-b),\quad a\in [-\frac12,\frac12]^N,\quad b\in[-\frac32,\frac32]^N.$$
If $Q\in\mathbb{D}$ is such that $3Q\not\subset\Omega,$ then the definition of $K_Q$ is the same, $a$ also runs through $[-\frac12,\frac12]^N,$ but the parallelepiped $b$ runs through is slightly smaller. {\color{red} to co-authors: if you wish, specify the domain of $b$ also.} {\it In what follows, we lighten the notations by pretending that $b$ is always defined on $[-\frac32,\frac32]^N.$}

It follows from the assumption on $K$ that
$$\|K_Q\|_{C^{2N+2}([-\frac32,\frac32]^N\times[-\frac32,\frac32]^N)}\leq\frac{c_{K,\theta}'\cdot\ell(Q)}{{\bf w}(Q)}.$$
Hence, there exists an extension $K_Q:\mathbb{T}^N\times\mathbb{T}^N\to\mathbb{C}$ such that
$$\|K_Q\|_{C^{2N+2}(\mathbb{T}^N\times\mathbb{T}^N)}\leq\frac{c_{K,\theta}''\cdot\ell(Q)}{{\bf w}(Q)}.$$
In particular,
$$\|K_Q\|_{W^{2N+2.2}(\mathbb{T}^N\times\mathbb{T}^N)}\leq \frac{c_{K,\theta}\cdot\ell(Q)}{{\bf w}(Q)}.$$
We may, therefore, write
$$K_Q(a,b)=\sum_{k,l\in\mathbb{Z}^N}{\rm coef}_{k,l,Q}e^{i\langle k,a\rangle}e^{i\langle l,b\rangle},\quad a,b\in\mathbb{T}^N,$$
where
$$\sum_{k,l\in\mathbb{Z}^N}(1+|k|^2+|l|^2)^{2N+2}|{\rm coef}_{k,l,Q}|^2\leq \Big(\frac{c_{K,\theta}\cdot\ell(Q)}{{\bf w}(Q)}\Big)^2.$$
In particular,
$$(1+|k|^2+|l|^2)^{N+1}\cdot|{\rm coef}_{k,l,Q}|\leq \frac{c_{K,\theta}\cdot\ell(Q)}{{\bf w}(Q)},\quad k,l\in\mathbb{Z}^N.$$

We now set
$$e_{k,l,Q}(x)=f(x)\chi_Q(x)e^{i\langle k,\frac{x-c_Q}{\ell(Q)}\rangle},$$
$$f_{k,l,Q}(y)=(1+|k|^2+|l|^2)^{N+1}{\rm coef}_{k,l,Q}\cdot \chi_{3Q}(y)e^{i\langle l,\frac{y-c_Q}{\ell(Q)}\rangle}.$$
This allows us to write
$$f(x)K(x,y)=\sum_{k,l\in\mathbb{Z}^N}(1+|k|^2+|l|^2)^{-N-1}\sum_{\substack{Q\in\mathbb{D}\\ Q\subset\Omega}}e_{k,l,Q}(x)f_{k,l,Q}(y).$$

Obviously,
$${\rm supp}(e_{k,l,Q})\subset Q,\quad {\rm supp}(f_{k,l,Q})\subset 3Q,$$
$$\|e_{k,l,Q}\|_{L_r(\mathbb{R}^N,{\bf w}(x)dx)}=\|f\|_{L_r(Q,{\bf w}(x)dx)},$$
$$\|f_{k,l,Q}\|_{L_{\infty}(\mathbb{R}^N)}=(1+|k|^2+|l|^2)^{N+1}|{\rm coef}_{k,l,Q}|\leq \frac{c_{K,\theta}\cdot\ell(Q)}{{\bf w}(Q)}.$$
\end{proof}


\begin{lemma}\label{RS intermediate estimate} Let $r>2.$ In the setting of Theorem \ref{RS-type theorem}, we have
$$\|M_fV\|_{\mathcal{L}_{N,\infty}(L_2(\Omega,{\bf w}(x)dx))}\leq c_{K,{\bf w},r}\Big\|\Big\{m(Q)^{\frac1N-\frac1r}\|f\|_{L_r(Q)}\Big\}_{\substack{Q\in\mathbb{D}\\ Q\subset\Omega}}
\Big\|_{N,\infty}.$$
\end{lemma}
\begin{proof} Let $A_{k,l}$ be the integral operator with integral kernel given by the formula
$$(x,y)\to\sum_{\substack{Q\in\mathbb{D}\\ Q\subset\Omega}}e_{k,l,Q}(x)f_{k,l,Q}(y),\quad x,y\in\Omega,\quad k,l\in\mathbb{Z}^N.$$
By Lemma \ref{lem:Weyl-off-diagonal-kernel-NWO}, we have
$$M_fV=\sum_{k,l\in\mathbb{Z}^N}(1+|k|^2+|l|^2)^{-N-1}A_{k,l}.$$
By the triangle inequality,
$$\|M_fV\|_{\mathcal{L}_{N,\infty}(L_2(\Omega,w_{\theta}(x)dx))}\leq$$
$$\leq \frac{N}{N-1}\sum_{k,l\in\mathbb{Z}^N}(1+|k|^2+|l|^2)^{-N-1} \|A_{k,l}\|_{\mathcal{L}_{N,\infty}(L_2(\Omega,w_{\theta}(x)dx))}.$$

By Lemma \ref{lem:Weyl-off-diagonal-kernel-NWO}, the sequences
$$\Big\{{\bf w}(Q)^{\frac1r-\frac12}\frac{e_{k,l,Q}}{\|f\|_{L_r(Q,{\bf w}(x)dx)}}\Big\}_{\substack{Q\in\mathbb{D}\\ Q\subset\Omega}},\quad \Big\{\frac{c_{K,\theta}^{-1}{\bf w}(Q)^{\frac12}}{\ell(Q)}f_{k,l,Q}\Big\}_{\substack{Q\in\mathbb{D}\\ Q\subset\Omega}}$$
are $(r,1)$-NWO and, respectively, $(\infty,3)$-NWO in the sense of Definition \ref{NWO def}. Using Lemma \ref{NWO lemma}, we conclude that
$$\|A_{k,l}\|_{\mathcal{L}_{N,\infty}(L_2(\Omega,{\bf w}(x)dx))}\leq$$
$$\leq c_{\theta,r,1,\infty,3,N}c_{K,\theta}\Big\|\Big\{\ell(Q){\bf w}(Q)^{-\frac1r}\|f\|_{L_r(Q,{\bf w}(x)dx)}\Big\}_{\substack{Q\in\mathbb{D}\\ Q\subset\Omega}}\Big\|_{N,\infty}.$$
Therefore,
$$\|M_fV\|_{\mathcal{L}_{N,\infty}(L_2(\mathbb{R}^N_+,{\bf w}(x)dx))}\leq$$
$$\leq \frac{N}{N-1}c_{\theta,r,1,\infty,3,N}c_{K,\theta}\sum_{k,l\in\mathbb{Z}^N}(1+|k|^2+|l|^2)^{-N-1}\times$$
$$\times \Big\|\Big\{\ell(Q){\bf w}(Q)^{-\frac1r}\|f\|_{L_r(Q,{\bf w}(x)dx)}\Big\}_{\substack{Q\in\mathbb{D}\\ Q\subset\Omega}}\Big\|_{N,\infty}.$$
Obviously,
$${\bf w}(Q)^{-\frac1r}\|f\|_{L_r(Q,{\bf w}(x)dx)}\leq \Big(\frac1{{\bf w}(Q)}\cdot\sup_{x\in Q}{\bf w}(x)\Big)^{\frac1r}\times \|f\|_{L_r(Q)},\quad Q\in\mathbb{D},\quad Q\subset\Omega.$$
Combining the last two estimates, we complete the proof.
\end{proof}

\begin{proof}[Proof of Theorem \ref{RS-type theorem}] Take $r\in(2,N).$ The assertion follows by combining Lemmas \ref{RS intermediate estimate} and \ref{lem:Cwikel-dyadic-distribution}.
\end{proof}


\section{Proof of special Cwikel estimate in Dunkl setting}

\begin{lemma}\label{heat kernel estimate} For every pair of multi-indices $\alpha,\beta\in\mathbb{Z}^N_+,$ for every $t>0,$ we have
$$|(\partial_x^{\alpha}\partial_y^{\beta}h_t)(x,y)|\leq c_{\kappa,\alpha,\beta}t^{1-\frac{|\alpha|_1+|\beta|_1}{2}}d(x,y)^{-2}w(B(x,t^{\frac12}))^{-1}\exp(-\frac{c_{\kappa}}{t}d(x,y)^2).$$
\end{lemma}
\begin{proof} This is (slightly weakened form of) Theorem 3.1 in \cite{MR4042416}.
\end{proof}

\begin{lemma}\label{zhijie kernel estimate} Suppose $s+N>2.$ We have
$$\int_0^{\infty}t^{-\frac{s}{2}}w(B(x,t^{\frac12}))^{-1}\exp(-\frac{c_{\kappa}}{t}d(x,y)^2)dt\leq c_{\kappa,s}w(B(x,d(x,y)))^{-1}\times d(x,y)^{2-s}.$$
\end{lemma}
\begin{proof} By \eqref{doublinginequality},
$$w(B(x,t^{\frac12}))^{-1}\leq c_{\kappa}^{(1)}w(B(x,d(x,y)))^{-1}\times
\begin{cases}
\left(\frac{d(x,y)}{t^{\frac12}}\right)^{\mathbf{N}},
&0<t\leq d(x,y)^2,\\
\left(\frac{d(x,y)}{t^{\frac12}}\right)^N,&t\geq d(x,y)^2.
\end{cases}$$
Therefore,
$$\int_0^{\infty}t^{-\frac{s}{2}}w(B(x,t^{\frac12}))^{-1}\exp(-\frac{c_{\kappa}}{t}d(x,y)^2)dt\leq$$
$$\leq c_{\kappa}^{(1)}w(B(x,d(x,y))^{-1}\times\int_0^{d(x,y)^2}t^{-\frac{s}{2}}\left(\frac{d(x,y)}{t^{\frac12}}\right)^{\mathbf{N}}\exp(-\frac{c_{\kappa}}{t}d(x,y)^2)dt+$$
$$+c_{\kappa}^{(1)}w(B(x,d(x,y)))^{-1}\times\int_{d(x,y)^2}^{\infty}t^{-\frac{s}{2}}\left(\frac{d(x,y)}{t^{\frac12}}\right)^Ndt=$$
$$=c_{\kappa}^{(1)}w(B(x,d(x,y)))^{-1}\times d(x,y)^{2-s}\times\int_0^1t^{-\frac{s+\mathbf{N}}{2}}\exp(-\frac{c_{\kappa}}{t})dt+$$
$$+c_{\kappa}^{(1)}w(B(x,d(x,y)))^{-1}\times d(x,y)^{2-s}\times\int_1^{\infty}t^{-\frac{s+N}{2}}dt.$$
\end{proof}

The following lemma establishes derivative estimates for the integral kernels of Dunkl-Riesz transforms.

\begin{lemma}\label{lem:Dunkl-Riesz-endpoint-derivatives} For every pair of multi-indices $\alpha,\beta,$ for every $1\leq j\leq N$ with $N+|\alpha|_1+|\beta|_1>3,$ we have 
$$|(\partial_x^{\alpha}\partial_y^{\beta}K_{\Delta_{{\rm Dunkl}}^{-\frac12}})(x,y)|
\leq c_{\kappa,\alpha,\beta}\frac{d(x,y)^{1-|\alpha|_1-|\beta|_1}}{w(B(x,d(x,y)))},\quad x,y\in\mathbb{R}^N.$$
{\color{red} an issue here is that for $\alpha=\beta=0,$ this imposes the condition $N\geq 4.$ it would be good to prove something like this under the assumption $N\geq 3.$}
\end{lemma}
\begin{proof} We have
$$\Delta_{{\rm Dunkl}}^{-\frac12}=\pi^{-\frac12}\int_0^{\infty}e^{-t\Delta_{{\rm Dunkl}}}\frac{dt}{t^{\frac12}}.$$
Hence,
$$K_{\Delta_{{\rm Dunkl}}^{-\frac12}}(x,y)=\pi^{-\frac12}\int_0^{\infty}h_t(x,y)\frac{dt}{t^{\frac12}}.$$

Using Lemma \ref{heat kernel estimate}, we estimate
$$|(\partial_x^{\alpha}\partial_y^{\beta}K_{R_{{\rm Dunkl},j}})(x,y)|\leq$$
$$\leq \frac{c_{\kappa,\alpha,\beta}}{\pi^{\frac12} d(x,y)^2}\int_0^{\infty}t^{-\frac{|\alpha|_1+|\beta|_1-1}{2}}w(B(x,t^{\frac12}))^{-1}\exp(-\frac{c_{\kappa}}{t}d(x,y)^2)dt.$$
Applying Lemma \ref{zhijie kernel estimate}, we conclude
$$|(\partial_x^{\alpha}\partial_y^{\beta}K_{\Delta_{{\rm Dunkl}}^{-\frac12}})(x,y)|\leq$$
$$\leq \pi^{-\frac12}c_{\kappa,\alpha,\beta}\times c_{\kappa,|\alpha|_1+|\beta|_1-1}w(B(x,d(x,y)))^{-1}\times d(x,y)^{1-|\alpha|_1-|\beta|_1}.$$
The assertion follows immediately.
\end{proof}

\begin{lemma}\label{supremum vs average} For every $Q\in\mathbb{D}$ with $Q\subset\Omega,$ we have
$$\sup_{x\in Q}w_{\theta}(x)\leq c_{\theta}\frac{w_{\theta}(Q)}{m(Q)}.$$
\end{lemma}
\begin{proof} We have
$$w(Q)\geq c_{\theta}\prod_{A\subset\{1,\cdots,N\}}(\ell_Q+\sum_{k\in A}(c_Q)_k)^{\theta_A}.$$
On the other hand, if $x\in Q,$ then $|x-c_Q|_{\infty}\leq \frac12\ell(Q)$ and, therefore,
$$\sum_{k\in A}x_k\leq \frac{|A|}{2}\ell_Q+\sum_{k\in A}(c_Q)_k\leq N(\ell_Q+\sum_{k\in A}(c_Q)_k).$$
Thus,
$$w(x)=\prod_{A\subset\{1,\cdots,N\}}(\sum_{k\in A}x_k)^{\theta_A}\leq N^{\sum_{A\subset\{1,\cdots,N\}}\theta_A}\prod_{A\subset\{1,\cdots,N\}}(\ell_Q+\sum_{k\in A}(c_Q)_k)^{\theta_A}.$$
This completes the proof.
\end{proof}


\begin{theorem}\label{dunkl cwikel theorem} Let $N\geq 4.$ We have
$$\|M_f\Delta_{{\rm Dunkl}}^{-\frac12}\Big\|_{\mathcal{L}_{\frac{N}{s},\infty}(L_2(\mathbb{R}^N,w(x)dx))}\leq c_{\kappa}\|f\|_{L_N(\mathbb{R}^N)}.$$
\end{theorem}
\begin{proof} It suffices to estimate, for every chamber $\Gamma$ and for every $g\in G,$ the operator $M_{\chi_{\Gamma}\cdot f}\Delta_{{\rm Dunkl}}^{-\frac12}M_{\chi_{g\Gamma}}.$ In other words, it suffices to estimate an operator with integral kernel
$$f(x)K_{{\rm Dunkl}^{-\frac12}}(x,y),\quad x,y\in\Gamma,\quad y\in g\Gamma,$$
on the Hilbert space $L_2(\mathbb{R}^N,w(x)dx).$ In other words, it suffices to estimate an operator with integral kernel
$$f(x)K_{{\rm Dunkl}^{-\frac12}}(x,gy),\quad x,y\in\Gamma,$$
on the Hilbert space $L_2(\Gamma,w(x)dx).$ 

Let $N_1={\rm dim}({\rm span}(R)),$ $N_2=N-N_1$ and let $\Omega=\mathbb{R}_+^{N_1}\times\mathbb{R}^{N_2}.$ For a concrete chamber $\Gamma,$ let $\Phi:\mathbb{R}^N\to\mathbb{R}^N$ be a linear isomorphism given in Theorem \ref{phi theorem} so that $\Phi(\Gamma)=\Omega.$ 

It suffices to estimate an operator with integral kernel
$$f(\Phi^{-1}x)K_{{\rm Dunkl}^{-\frac12}}(\Phi^{-1}x,g\Phi^{-1}y),\quad x,y\in\Omega,$$
on the Hilbert space $L_2(\Omega,w(\Phi^{-1}x)dx).$ Using Lemma \ref{s2 trivial lemma}, it suffices to estimate an operator with integral kernel
$$f(\Phi^{-1}x)K_{{\rm Dunkl}^{-\frac12}}(\Phi^{-1}x,g\Phi^{-1}y),\quad x,y\in\Omega,$$
on the Hilbert space $L_2(\Omega,w_{\theta}(x)dx).$

By Lemma \ref{lem:Dunkl-Riesz-endpoint-derivatives},
$$\Big|\partial^{\alpha}_x\partial^{\beta}_yK_{{\rm Dunkl}^{-\frac12}}(\Phi^{-1}x,g\Phi^{-1}y)\Big|\leq c_{\alpha,\beta,\kappa}\frac{|x-y|^{1-|\alpha|_1-|\beta|_1}}{w_{\theta}(B(x,|x-y|))},\quad x,y\in\Omega.$$
The other assumptions in Theorem \ref{RS-type theorem} follow from Lemma \ref{supremum vs average}, Lemma \ref{s2 upper estimate} and Lemma \ref{s2 lower estimate}. The assertion follows now from Theorem \ref{RS-type theorem}.
\end{proof}

\section{Proof of Theorem \ref{main theorem} \eqref{mta}}

\begin{lemma}\label{tj commutator fact} We have
$$[T_{e_j},M_f]=M_{\partial_jf}+\frac12\sum_{v\in R}v_j\kappa(v)M_{\delta_vf}O_v,$$
where 
$$(\delta_vf)(x)=\frac{f(x)-(O_vf)(x)}{\langle v,x\rangle},\quad v\in R,\quad x\in\mathbb{R}^N,$$
and where $O_v$ is (the operation of composition with) reflection with respect to vector $v.$
\end{lemma}
\begin{proof} Obvious.
\end{proof}

\begin{lemma}\label{deltav estimate} We have
$$\|\delta_vf\|_{L_N(\mathbb{R}^N)}\leq c_N\|f\|_{\dot{W}^{1,N}(\mathbb{R}^N)}.$$
\end{lemma}
\begin{proof} Without loss of generality, $v=e_N.$ Set
$$(D_th)(x)=h(x_1,\cdots,x_{N-1},tx_N).$$
We have
$$\delta_vf=\int_0^1D_t(\partial_Nf)dt.$$
Thus,
$$\|\delta_vf\|_{L_N(\mathbb{R}^N)}\leq\int_0^1\|D_t(\partial_Nf)\|_{L_N(\mathbb{R}^N)}dt=$$
$$=\int_0^1t^{-\frac1N}\|\partial_Nf\|_{L_N(\mathbb{R}^N)}dt=\frac{N}{N-1}\|\partial_Nf\|_{L_N(\mathbb{R}^N)}.$$
\end{proof}

\begin{lemma}\label{from doi lemma} For every $N\geq1,$ we have
$$\|[B,A]B^{-1}\|_{N,\infty}\leq c_N\|B^{-1}[B^2,A]B^{-1}\|_{N,\infty}.$$
\end{lemma}

\begin{proof}[Proof of Theorem \ref{main theorem} \eqref{mta}] We prove the assertion under {\it a priori} restriction $f\in C^{\infty}_c(\mathbb{R}^N).$ Passing to the case of $f\in\dot{W}^{1,N}(\mathbb{R}^N)$ is standard (see e.g {\color{red} refer to our papers}).

We have
$$[R_{{\rm Dunkl},j},M_f]=[T_{e_j}\Delta_{{\rm Dunkl}}^{-\frac12},M_f]=[T_{e_j},M_f]\Delta_{{\rm Dunkl}}^{-\frac12}+T_{e_j}[\Delta_{{\rm Dunkl}}^{-\frac12},M_f]=$$
$$=[T_{e_j},M_f]\Delta_{{\rm Dunkl}}^{-\frac12}-R_{{\rm Dunkl},j}[\Delta_{{\rm Dunkl}}^{\frac12},M_f]\Delta_{{\rm Dunkl}}^{-\frac12}.$$
By triangle inequality,
$$\|[R_{{\rm Dunkl},j},M_f]\|_{\mathcal{L}_{N,\infty}(L_2(\mathbb{R}^N,w(x)dx))}\leq$$
$$\leq\frac{N}{N-1}\|[T_{e_j},M_f]\Delta_{{\rm Dunkl}}^{-\frac12}\|_{\mathcal{L}_{N,\infty}(L_2(\mathbb{R}^N,w(x)dx))}+$$
$$+\frac{N}{N-1}\|[\Delta_{{\rm Dunkl}}^{\frac12},M_f]\Delta_{{\rm Dunkl}}^{-\frac12}\|_{\mathcal{L}_{N,\infty}(L_2(\mathbb{R}^N,w(x)dx))}.$$
Using Lemma \ref{from doi lemma}, we write
$$\|[R_{{\rm Dunkl},j},M_f]\|_{\mathcal{L}_{N,\infty}(L_2(\mathbb{R}^N,w(x)dx))}\leq$$
$$\leq\frac{N}{N-1}\|[T_{e_j},M_f]\Delta_{{\rm Dunkl}}^{-\frac12}\|_{\mathcal{L}_{N,\infty}(L_2(\mathbb{R}^N,w(x)dx))}+$$
$$+\frac{N}{N-1}c_N\|\Delta_{{\rm Dunkl}}^{-\frac12}[\Delta_{{\rm Dunkl}},M_f]\Delta_{{\rm Dunkl}}^{-\frac12}\|_{\mathcal{L}_{N,\infty}(L_2(\mathbb{R}^N,w(x)dx))}.$$

Note that
$$\Delta_{{\rm Dunkl}}^{-\frac12}[\Delta_{{\rm Dunkl}},M_f]\Delta_{{\rm Dunkl}}^{-\frac12}=$$
$$=\sum_{k=1}^N\Delta_{{\rm Dunkl}}^{-\frac12}[T_{e_k},M_f]R_{{\rm Dunkl},k}+R_{{\rm Dunkl},k}[T_{e_k},M_f]\Delta_{{\rm Dunkl}}^{-\frac12}.$$
By triangle inequality,
$$\|\Delta_{{\rm Dunkl}}^{-\frac12}[\Delta_{{\rm Dunkl}},M_f]\Delta_{{\rm Dunkl}}^{-\frac12}\|_{\mathcal{L}_{N,\infty}(L_2(\mathbb{R}^N,w(x)dx))}\leq$$
$$\leq\frac{N}{N-1}\sum_{k=1}^N\|\Delta_{{\rm Dunkl}}^{-\frac12}[T_{e_k},M_f]\|_{\mathcal{L}_{N,\infty}(L_2(\mathbb{R}^N,w(x)dx))}+$$
$$+\frac{N}{N-1}\sum_{k=1}^N\|[T_{e_k},M_f]\Delta_{{\rm Dunkl}}^{-\frac12}\|_{\mathcal{L}_{N,\infty}(L_2(\mathbb{R}^N,w(x)dx))}.$$

Combining the estimates in the last two paragraphs, we obtain
$$\|[R_{{\rm Dunkl},j},M_f]\|_{\mathcal{L}_{N,\infty}(L_2(\mathbb{R}^N,w(x)dx))}\leq$$
$$\leq c_N'\sum_{k=1}^N\|\Delta_{{\rm Dunkl}}^{-\frac12}[T_{e_k},M_f]\|_{\mathcal{L}_{N,\infty}(L_2(\mathbb{R}^N,w(x)dx))}+$$
$$+c_N'\sum_{k=1}^N\|[T_{e_k},M_f]\Delta_{{\rm Dunkl}}^{-\frac12}\|_{\mathcal{L}_{N,\infty}(L_2(\mathbb{R}^N,w(x)dx))}.$$
Using Lemma \ref{tj commutator fact}, the fact that $\Delta_{{\rm Dunkl}}$ commutes with $O_v$ and the triangle inequality, we obtain
$$\|[T_{e_k},M_f]\Delta_{{\rm Dunkl}}^{-\frac12}\|_{\mathcal{L}_{N,\infty}(L_2(\mathbb{R}^N,w(x)dx))}\leq$$
$$\leq\frac{N}{N-1}\|M_{\partial_kf}\Delta_{{\rm Dunkl}}^{-\frac12}\|_{\mathcal{L}_{N,\infty}(L_2(\mathbb{R}^N,w(x)dx))}+$$
$$+\frac{N}{2(N-1)}\sum_{v\in R}|v|\kappa(v)\|M_{\delta_vf}\Delta_{{\rm Dunkl}}^{-\frac12}\|_{\mathcal{L}_{N,\infty}(L_2(\mathbb{R}^N,w(x)dx))}.$$
Thus,
$$\|[R_{{\rm Dunkl},j},M_f]\|_{\mathcal{L}_{N,\infty}(L_2(\mathbb{R}^N,w(x)dx))}\leq$$
$$\leq c_N''\sum_{k=1}^N\|M_{\partial_kf}\Delta_{{\rm Dunkl}}^{-\frac12}\|_{\mathcal{L}_{N,\infty}(L_2(\mathbb{R}^N,w(x)dx))}+$$
$$+c_{\kappa}\sum_{v\in R}\|M_{\delta_vf}\Delta_{{\rm Dunkl}}^{-\frac12}\|_{\mathcal{L}_{N,\infty}(L_2(\mathbb{R}^N,w(x)dx))}.$$
The assertion follows now from Theorem \ref{dunkl cwikel theorem} and Lemma \ref{deltav estimate}.
\end{proof}

\section{Proof of Theorem \ref{main theorem} \eqref{mtb}}\label{lower estimate section}

\begin{lemma}\label{nondegenerate} Let $1\leq j\leq N.$ If $x_j\geq y_j,$ then
$$K_{{\rm Dunkl},j}(x,y)\geq c_{R,k}\frac{x_j-y_j}{|x-y|}\cdot\frac1{w(B(x,|x-y|))}.$$
If $x_j\leq y_j,$ then
$$K_{{\rm Dunkl},j}(x,y)\leq c_{R,k}\frac{x_j-y_j}{|x-y|}\cdot\frac1{w(B(x,|x-y|))}.$$
\end{lemma}
\begin{proof} The proof is identical to that in \cite[Theorem 1.2]{MR4733961}.
\end{proof}


\begin{lemma}\label{nondegenerate corollary} Let $Q\in\mathbb{D}.$ We have
$$K_j(x,y)\geq \frac{c_k}{w(Q)},\quad x\in 3Q,\quad y\in Q+3s(Q)e_j.$$
\end{lemma}
\begin{proof} The assertion follows immediately from Lemma \ref{nondegenerate}.
\end{proof}

\begin{lemma}\label{improved zhijie lemma} Let real-valued $f\in L_1^{{\rm loc}}(\mathbb{R}^N)$ be such that the commutator $[R_{{\rm Dunkl},j},M_f]$ is bounded on $L_2(\mathbb{R}^N,w(x)dx).$ For every $Q\in\mathbb{D},$ there exist $\xi_Q$ and $\eta_Q$ supported in $7Q$ such that $\|\xi_Q\|_{\infty},\|\eta_Q\|_{\infty}\leq m(Q)^{-\frac12}$ and such that
$$\frac1{w(3Q)}\int_{3Q}|f(y)-f_{3Q}|w(y)dy\leq c_k|\langle[R_{{\rm Dunkl},j},M_f]\xi_Q,\eta_Q\rangle_{L_2(\mathbb{R}^N,w(x)dx)}|.$$
\end{lemma}
\begin{proof} Let $\alpha$ be the median of $f$ on $3Q$ and $\beta$ be the median of $f$ on $Q+3s(Q)e_j.$ Assume for definiteness $\alpha\geq\beta.$ Set
$$A_1=\{x\in 3Q:\ f(x)\geq\alpha\},\quad B_1=\{x\in Q+3s(Q)e_j:\ f(x)\leq\beta\},$$
$$A_2=\{x\in 3Q:\ f(x)\leq\beta\},\quad B_2=\{x\in Q+3s(Q)e_j:\ f(x)\geq\beta\},$$
$$\xi_{Q,s}=w(Q)^{-\frac12}\chi_{A_s},\quad \eta_{Q,s}=w(Q)^{-\frac12}\chi_{B_s},\quad s=1,2.$$

We have
$$\langle[R_{{\rm Dunkl},j},M_f]\xi_{Q,s},\eta_{Q,s}\rangle_{L_2(\mathbb{R}^N,w(x)dx)}=$$
$$=\frac1{w(Q)}\int_{B_s\times A_s}K_j(x,y)(f(y)-f(x))w(x)w(y)dxdy.$$

If $x\in B_1$ and $y\in A_1,$ then
$$f(y)-f(x)=(f(y)-\alpha)+(\alpha-\beta)+(\beta-f(x))=$$
$$=|f(y)-\alpha|+|\alpha-\beta|+|f(x)-\beta|\geq |f(y)-\alpha|+|\alpha-\beta|\geq0.$$
If $x\in B_2$ and $y\in A_2,$ then
$$f(y)-f(x)=(f(y)-\beta)+(\beta-f(x))=$$
$$=-|f(y)-\beta|-|f(x)-\beta|\leq -|f(y)-\beta|\leq 0.$$

Thus,
$$\langle[R_{{\rm Dunkl},j},M_f]\xi_{Q,1},\eta_{Q,1}\rangle_{L_2(\mathbb{R}^N,w(x)dx)}\geq $$
$$\geq \frac1{w(Q)}\int_{B_1\times A_1}K_j(x,y)\big(|f(y)-\alpha|+|\alpha-\beta|\big)w(x)w(y)dxdy\geq $$
$$\geq\frac{c_k}{w^2(Q)}\int_{B_1\times A_1}\big(|f(y)-\alpha|+|\alpha-\beta|\big)w(x)w(y)dxdy=$$
$$=\frac{c_k}{w^2(Q)}\cdot\Big(w(B_1)\cdot\int_{A_1}|f(y)-\alpha|w(y)dy+|\alpha-\beta|w(B_1)w(A_1)\Big)\geq $$
$$=\frac{c_k}{w^2(Q)}\cdot\Big(\frac12w(Q)\cdot\int_{A_1}|f(y)-\alpha|w(y)dy+|\alpha-\beta|\cdot \frac14w^2(Q)\Big)\geq $$
$$\geq\frac{c_k}{4w(Q)}\cdot\Big(\int_{A_1}|f(y)-\alpha|w(y)dy+|\alpha-\beta|w(Q)\Big)$$
and
$$-\langle[R_{{\rm Dunkl},j},M_f]\xi_{Q,2},\eta_{Q,2}\rangle_{L_2(\mathbb{R}^N,w(x)dx)}\geq $$
$$\geq \frac1{w(Q)}\int_{B_2\times A_2}K_j(x,y)|f(y)-\beta|w(x)w(y)dxdy\geq $$
$$\geq\frac{c_k}{w^2(Q)}\int_{B_2\times A_2}|f(y)-\beta|w(x)w(y)dxdy=$$
$$=\frac{c_k}{w^2(Q)}\cdot w(B_2)\cdot\int_{A_2}|f(y)-\beta|w(y)dy\geq $$
$$=\frac{c_k}{2w(Q)}\int_{A_2}|f(y)-\beta|w(y)dy.$$

It remains to note that
$$\int_{3Q}|f(y)-\alpha|w(y)dy=\int_{A_1}|f(y)-\alpha|w(y)dy+$$
$$+\int_{A_2}|f(y)-\alpha|w(y)dy+\int_{3Q\backslash(A_1\cup A_2)}|f(y)-\alpha|w(y)dy\leq$$
$$\leq\int_{A_1}|f(y)-\alpha|w(y)dy+\int_{A_2}|f(y)-\beta|w(y)dy$$
$$+\int_{A_2}|\alpha-\beta|w(y)dy+\int_{3Q\backslash(A_1\cup A_2)}|\alpha-\beta|w(y)dy\leq$$
$$\leq\int_{A_1}|f(y)-\alpha|w(y)dy+\int_{A_2}|f(y)-\beta|w(y)dy+2|\alpha-\beta|w(3Q).$$
Combining the last $3$ equations, we conclude
$$\frac{c_k'}{w(Q)}\int_{3Q}|f(y)-\alpha|w(y)dy\leq\max_{s=1,2}|\langle[R_{{\rm Dunkl},j},M_f]\xi_{Q,s},\eta_{Q,s}\rangle_{L_2(\mathbb{R}^N,w(x)dx)}|.$$

If
$$|\langle[R_{{\rm Dunkl},j},M_f]\xi_{Q,1},\eta_{Q,1}\rangle_{L_2(\mathbb{R}^N,w(x)dx)}|\geq |\langle[R_{{\rm Dunkl},j},M_f]\xi_{Q,2},\eta_{Q,2}\rangle_{L_2(\mathbb{R}^N,w(x)dx)}|,$$
then we set $\xi_Q=\xi_{Q,1}$ and $\eta=\eta_{Q,1}.$ Otherwise, set $\xi_Q=\xi_{Q,2}$ and $\eta_Q=\eta_{Q,2}.$ Thus,
$$\frac{c_k'}{w(Q)}\int_{3Q}|f(y)-\alpha|w(y)dy\leq|\langle[R_{{\rm Dunkl},j},M_f]\xi_Q,\eta_Q\rangle_{L_2(\mathbb{R}^N,w(x)dx)}|.$$
Hence,
$$c_k'|f_{3Q}-\alpha|\leq|\langle[R_{{\rm Dunkl},j},M_f]\xi_Q,\eta_Q\rangle_{L_2(\mathbb{R}^N,w(x)dx)}|$$
and, finally,
$$\frac{c_k'}{w(Q)}\int_{3Q}|f(y)-f_{3Q}|w(y)dy\leq 2|\langle[R_{{\rm Dunkl},j},M_f]\xi_Q,\eta_Q\rangle_{L_2(\mathbb{R}^N,w(x)dx)}|.$$
\end{proof}

\begin{lemma}\label{standard} Suppose that $f\in L_1^{{\rm loc}}(\mathbb{R}^N).$ If $[R_{{\rm Dunkl},j},M_f]\in\mathcal{L}_{N,\infty}(L_2(\mathbb{R}^N,w(x)dx))$ for some $j\in\{1,2,\cdots,N\},$ then
$$\Big\|\Big\{\frac1{w(3Q)}\int_{3Q}|f(x)-f_{3Q}|w(x)dx\Big\}_{Q\in\mathbb{D}}\Big\|_{l_{N,\infty}(\mathbb{D})}\leq$$
$$\leq c_k\|[R_{{\rm Dunkl},j},M_f]\|_{\mathcal{L}_{N,\infty}(L_2(\mathbb{R}^N,w(x)dx))}.$$
\end{lemma}
\begin{proof} Let $\{\xi_Q\}_{Q\in\mathbb{D}}$ and $\{\eta_Q\}_{Q\in\mathbb{D}}$ be functions constructed in Lemma \ref{improved zhijie lemma}. By Lemma \ref{NWO lower estimate},
$$\Big\{\langle [R_{{\rm Dunkl},j},M_f]\xi_Q,\eta_Q\rangle_{L_2(\mathbb{R}^N,w(x)dx)}\Big\}_{Q\in\mathbb{D}}\in l_p.$$
The first assertion follows now from Lemma \ref{improved zhijie lemma}.
\end{proof}

\begin{lemma}\label{for reverse holder first lemma} Let $Q$ be a cube in $\mathbb{R}^N$ (with sides parallel to coordinate hyperplanes). We have
$$\sup_{x\in Q}w(x)\leq \frac{c_k}{m(Q)}\int_Qw(x)dx.$$
\end{lemma}
\begin{proof} If $x=c_Q+y,$ $y\in [-\frac12s(Q),\frac12s(Q)]^d,$ then
$$|\langle x,v\rangle|\leq |\langle c_Q,v\rangle|+r|\langle y,v\rangle|\leq |\langle c_Q,v\rangle|+\frac{\sqrt{d}}{2}s(Q).$$
Thus,
$$w(x)=\prod_{v\in R}|\langle x,v\rangle|^{k(v)}\leq\prod_{v\in R}(|\langle c_Q,v\rangle|+\frac{\sqrt{d}}{2}s(Q))^{k(v)}.$$
In particular,
$$LHS\leq \prod_{v\in R}(|\langle c_Q,v\rangle|+\frac{\sqrt{d}}{2}s(Q))^{k(v)}.$$

Fix chamber $\Gamma$ such that $c_Q\in\Gamma.$ We have
$$\int_Qw(x)dx\geq\int_{c_Q+(B(0,\frac12s(Q))\cap\Gamma)}w(x)dx=\int_{B(0,\frac12s(Q))\cap\Gamma}w(c_Q+y)dy.$$
If $y\in\Gamma,$ then
$$|\langle c_Q+y,v\rangle|=|\langle c_Q,v\rangle|+|\langle y,v\rangle|,\quad v\in R.$$
Thus,
$$w(c_Q+y)=\prod_{v\in R}(|\langle c_Q,v\rangle|+|\langle y,v\rangle|)^{k(v)}.$$
Hence,
$$\int_Qw(x)dx\geq\int_{B(0,\frac12s(Q))\cap\Gamma}\prod_{v\in R}(|\langle c_Q,v\rangle|+|\langle y,v\rangle|)^{k(v)}dy.$$
Fix a closed cone $\Gamma'\subset\Gamma$ such that (i) interior of $\Gamma'$ is non-empty and (ii) $\partial\Gamma\cap\partial\Gamma'=\{0\}.$ We have
$$|\langle y,v\rangle|\geq c_{\Gamma'}|y|,\quad y\in\Gamma',\quad v\in R.$$
We, therefore, have
$$\int_Qw(x)dx\geq\int_{B(0,\frac12s(Q))\cap\Gamma'}\prod_{v\in R}(|\langle c_Q,v\rangle|+c_{\Gamma'}|y|)^{k(v)}dy\geq$$
$$\geq \int_{\substack{y\in\Gamma'\\ \frac14s(Q)\leq|y|\leq \frac14s(Q)}}\prod_{v\in R}(|\langle c_Q,v\rangle|+\frac12c_{\Gamma'}s(Q))^{k(v)}dy=$$
$$=s(Q)^N\prod_{v\in R}(|\langle c_Q,v\rangle|+\frac14c_{\Gamma'}s(Q))^{k(v)}\times (2^{-N}-4^{-N})\int_{\substack{y\in\Gamma'\\ |y|\leq1}}dy.$$
Combining this with the preceding paragraph, we complete the proof.
\end{proof}

\begin{lemma}\label{for reverse holder second lemma} Suppose $1<t<1+\frac1{|R|\cdot\max_{v\in R}k(v)}.$ Let $Q$ be a cube in $\mathbb{R}^N$ (with sides parallel to coordinate hyperplanes). We have
$$\Big(\frac{\int_Qw^{1-t}(x)dx}{\int_Qw(x)dx}\Big)^{\frac1t}\leq c_{t,k} \frac{m(Q)}{\int_Qw(x)dx}.$$
\end{lemma}
\begin{proof} If $v\in R$ is such that
$$\sqrt{d}s(Q)\leq |\langle c_Q,v\rangle|,$$
then
$$|\langle x,v\rangle|\geq |\langle c_Q,v\rangle|-|\langle x-c_Q,v\rangle|\geq |\langle c_Q,v\rangle|-\frac{\sqrt{d}}{2}s(Q)\geq \frac12|\langle c_Q,v\rangle|,\quad x\in Q.$$
In this case,
$$\int_Q|\langle x,v\rangle|^{(1-t)|R|k(v)}dx\leq 2^{(t-1)|R|k(v)}|\langle c_Q,v\rangle|^{(1-t)|R|k(v)}\cdot\int_Qdx=$$
$$=2^{(t-1)|R|k(v)}|\langle x_0,v\rangle|^{(1-t)|R|k(v)}\cdot s(Q)^N.$$

If $v\in R$ is such that
$$|\langle c_Q,v\rangle|\leq \sqrt{d}s(Q),$$
then let $y_0=\langle c_Q,v\rangle$ and $z_0=c_Q-y_0v.$ Let us introduce new variables $y=\langle x,v\rangle$ and $z=x-yv$ (it varies in the plane orthogonal to $v$). We write
$$\int_Q|\langle x,v\rangle|^{(1-t)|R|k(v)}dx\leq\int_{|x-c_Q|\leq\frac{\sqrt{d}}{2}s(Q)}|\langle x,v\rangle|^{(1-t)|R|k(v)}dx\leq$$
$$\leq \int_{|y-y_0|,|z-z_0|\leq \frac{\sqrt{d}}{2}r}|y|^{(1-t)|R|k(v)}dydz=$$
$$=(\frac{\sqrt{d}}{2}s(Q))^{N-1}{\rm Vol}(\mathbb{B}^{N-1})\cdot\int_{|y-y_0|\leq \frac{\sqrt{d}}{2}s(Q)}|y|^{(1-t)|R|k(v)}dy\leq$$
$$=(\frac{\sqrt{d}}{2}s(Q))^{N-1}{\rm Vol}(\mathbb{B}^{N-1})\cdot\int_{|y|\leq \frac{3\sqrt{d}}{2}s(Q)}|y|^{(1-t)|R|k(v)}dy=$$
$$=s(Q)^{N+(1-t)|R|k(v)}\cdot (\frac{\sqrt{d}}{2})^{N-1}{\rm Vol}(\mathbb{B}^{N-1})\cdot\int_{|y|\leq \frac{3\sqrt{d}}{2}}|y|^{(1-t)|R|k(v)}dy.$$

Combining the estimates in the preceding paragraphs, we write
$$\int_Q|\langle x,v\rangle|^{(1-t)|R|k(v)}dx\leq c_1s(Q)^N\big(\max\{|\langle c_Q,v\rangle|,r\}\big)^{(1-t)|R|k(v)}.$$
By H\"older inequality,
$$\int_Qw^{1-t}(x)dx\leq \prod_{v\in R}\Big(\int_Q|\langle x,v\rangle|^{(1-t)|R|k(v)}dx\Big)^{\frac1{|R|}}\leq$$
$$\leq c_2s(Q)^N\prod_{v\in R}\big(\max\{|\langle c_Q,v\rangle|,s(Q)\}\big)^{(1-t)k(v)}.$$
The assertion follows easily.
\end{proof}

\begin{lemma}\label{hytonen verification lemma} Measures $dx$ and $w(x)dx$ (and distance $(x,y)\to |x-y|_{\infty}$) satisfy the conditions in Proposition 1.2 in \cite{Hytonen-new}.
\end{lemma}
\begin{proof} That both measures are doubling is obvious.

Let $Q$ be a cube in $\mathbb{R}^N$ (with sides parallel to coordinate hyperplanes). By Lemmas \ref{for reverse holder first lemma} and \ref{for reverse holder second lemma},
$$\Big(\frac1{m(Q)}\int_Qw^t(x)dx\Big)^{\frac1t}\leq\frac{c_k}{m(Q)}\int_Qw(x)dx,\quad t>0.$$
$$\Big(\frac{\int_Qw^{1-t}(x)dx}{\int_Qw(x)dx}\Big)^{\frac1t}\leq c_G \frac{m(Q)}{\int_Qw(x)dx},\quad 1<t<1+\frac1{|R|\cdot\max_{v\in R}k(v)}.$$
The assertion follows now from Proposition 3.6 in \cite{Hytonen-new}.
\end{proof}

\begin{proof}[Proof of Theorem \ref{main theorem} \eqref{mtb}] The left hand side in Lemma \ref{standard} is the oscillation semi-norm of the function $f$ with respect to measure $w(x)dx$ and distance $(x,y)\to |x-y|_{\infty}.$ So, the oscillation semi-norm of the function $f$ with respect to measure $w(x)dx$ and distance $(x,y)\to |x-y|_{\infty}$ is controlled by the $\|\cdot\|_{\mathcal{L}_{N,\infty}(L_2(\mathbb{R}^N,w(x)dx))}$-norm of the commutator $[R_{{\rm Dunkl},j},M_f].$

Lemma \ref{hytonen verification lemma} asserts that measures $dx$ and $w(x)dx$ (and distance $(x,y)\to |x-y|_{\infty}$) satisfy the conditions in Proposition 1.2 in \cite{Hytonen-new}. Appealing to Proposition 1.2 in \cite{Hytonen-new}, we conclude that the oscillation semi-norm of the function $f$ with respect to measure $dx$ and distance $(x,y)\to |x-y|_{\infty}$ is controlled by the one with respect to measure $w(x)dx$ and distance $(x,y)\to |x-y|_{\infty}.$ By the preceding paragraph, the oscillation semi-norm of the function $f$ with respect to measure $dx$ and distance $(x,y)\to |x-y|_{\infty}$ is controlled by the $\|\cdot\|_{\mathcal{L}_{N,\infty}(L_2(\mathbb{R}^N,w(x)dx))}$-norm of the commutator $[R_{{\rm Dunkl},j},M_f].$

That oscillation semi-norm of the function $f$ with respect to measure $dx$ and distance $(x,y)\to |x-y|_{\infty}$ is equivalent to the Sobolev semi-norm of the function $f$ is a standard fact. The proof (in much more general setting) is available in Proposition 1.29 in \cite{HKarxiv}. The assertion follows now from the preceding paragraph. 
\end{proof}

\section{Proof of smoothness}

\begin{theorem}\label{smoothness theorem} Let $f\in C^{\infty}_c(\mathbb{R}^N).$
\begin{enumerate}[{\rm (i)}]
\item\label{smootha} for every $n\in\mathbb{N},$ $s_1,s_2\in\mathbb{R}$ with $s_1+s_2=n,$ the operator\footnote{Commutators are repeated $n$ times}
$$(1+\Delta_{{\rm Dunkl}})^{-\frac{s_1}{2}}[\Delta_{{\rm Dunkl}},[\Delta_{{\rm Dunkl}},[\cdots,M_f]\cdots](1+\Delta_{{\rm Dunkl}})^{-\frac{s_2}{2}}$$
is bounded on $L_2(\mathbb{R}^N,w(x)dx)$;
\item\label{smoothb} for every $n\in\mathbb{N},$ $s_1,s_2\in\mathbb{R}$ with $s_1+s_2>n,$ the operator\footnote{Commutators are repeated $n$ times}
$$(1+\Delta_{{\rm Dunkl}})^{-\frac{s_1}{2}}[\Delta_{{\rm Dunkl}},[\Delta_{{\rm Dunkl}},[\cdots,M_f]\cdots](1+\Delta_{{\rm Dunkl}})^{-\frac{s_2}{2}}$$
belongs to $\mathcal{L}_{\frac{N}{s_1+s_2-n},\infty}(L_2(\mathbb{R}^N,w(x)dx));$
\end{enumerate}
\end{theorem}

\begin{lemma}\label{commutator with delta} Let $f\in C^{\infty}_c(\mathbb{R}^N).$ We have
$$[\Delta_{{\rm Dunkl}},M_f]=2\sum_{j=1}^NM_{\partial_jf}T_{e_j}+\frac12\sum_{j=1}^N\sum_{v\in R}v_j\kappa(v)M_{\delta_vf}(T_{e_j}O_v+O_vT_{e_j})+\sum_{g\in G}M_{f_g}U_g.$$
All functions on the right hand side are smooth and compactly supported.
\end{lemma}
\begin{proof} We have
$$[\Delta_{{\rm Dunkl}},M_f]=\sum_{j=1}^NT_{e_j}[T_{e_j},M_f]+[T_{e_j,M_f}T_{e_j}]=$$
$$=\sum_{j=1}^NT_{e_j}M_{\partial_jf}+M_{\partial_jf}T_{e_j}+\frac12\sum_{j=1}^N\sum_{v\in R}v_j\kappa(v)\big(T_{e_j}M_{\delta_vf}O_v+M_{\delta_vf}O_vT_{e_j}\big)=$$
$$=2\sum_{j=1}^NM_{\partial_jf}T_{e_j}+\frac12\sum_{j=1}^N\sum_{v\in R}v_j\kappa(v)\big(M_{\delta_vf}T_{e_j}O_v+M_{\delta_vf}O_vT_{e_j}\big)+$$
$$+\sum_{j=1}^N[T_{e_j},M_{\partial_jf}]+\frac12\sum_{j=1}^N\sum_{v\in R}v_j\kappa(v)[T_{e_j},M_{\delta_vf}]O_v=$$
$$=2\sum_{j=1}^NM_{\partial_jf}T_{e_j}+\frac12\sum_{j=1}^N\sum_{v\in R}v_j\kappa(v)\big(M_{\delta_vf}T_{e_j}O_v+M_{\delta_vf}O_vT_{e_j}\big)+$$
$$+\sum_{j=1}^NM_{\partial_j^2f}+\frac12\sum_{j=1}^N\sum_{v\in R}\alpha_j\kappa(v)M_{\delta_v(\partial_jf)}O_v+$$
$$+\frac12\sum_{j=1}^N\sum_{v\in R}v_j\kappa(v)M_{\partial_j(\delta_vf)}O_v+\frac14\sum_{j=1}^N\sum_{v,v'\in R}v_jv_j'\kappa(v)\kappa(v')M_{\delta_{v'}\delta_vf}O_{v'}O_v.$$
\end{proof}

\begin{lemma}\label{base of induction} For every $f\in C^{\infty}_c(\mathbb{R}^N),$ the operator
$$[\Delta_{{\rm Dunkl}},M_f](1+\Delta_{{\rm Dunkl}})^{-\frac12}$$
is bounded and
$$(1+\Delta_{{\rm Dunkl}})^{-1}[\Delta_{{\rm Dunkl}},M_f]\in \mathcal{L}_{N,\infty}(L_2(\mathbb{R}^N,w(x)dx)).$$
\end{lemma}
\begin{proof} We prove only the first assertion. The proof of the second one is similar (modulo Theorem \ref{dunkl cwikel theorem}).

Since each $T_{e_j}$ commutes  with $(1+\Delta_{{\rm Dunkl}})^{-\frac12},$ each $O_v$ commutes with $(1+\Delta_{{\rm Dunkl}})^{-\frac12}$ and since the operators
$$\frac{T_{e_j}}{(1+\Delta_{{\rm Dunkl}})^{\frac12}},\quad 1\leq j\leq N,$$
are bounded, the assertion follows from Lemma \ref{commutator with delta}.
\end{proof}

\begin{lemma}\label{step of induction} For every $n\in\mathbb{N},$ the operator\footnote{Commutators are repeated $n$ times}
$$[\Delta_{{\rm Dunkl}},[\Delta_{{\rm Dunkl}},[\cdots,M_f]\cdots](1+\Delta_{{\rm Dunkl}})^{-\frac{n}{2}}$$
is bounded.
\end{lemma}
\begin{proof} Denote for brevity
$${\rm Expr}_n(f)=[\Delta_{{\rm Dunkl}},[\Delta_{{\rm Dunkl}},[\cdots,M_f]\cdots](1+\Delta_{{\rm Dunkl}})^{-\frac{n}{2}}.$$

We prove the assertion by induction on $f.$ The $n$-th statement is "for every $f\in C^{\infty}_c(\mathbb{R}^N),$ the operator
${\rm Expr}_n(f)$ is bounded." Base of induction is given by Lemma \ref{base of induction}. It suffices to prove the step of induction. 

Since each $T_{e_j}$ and each $O_v$ commutes with $\Delta_{{\rm Dunkl}},$ it follows from Lemma \ref{commutator with delta} that
$${\rm Expr}_{n+1}(f)=$$
$$=\sum_{j=1}^N{\rm Expr}_n(f_j^{(1)})\cdot \frac{T_{e_j}}{(1+\Delta_{{\rm Dunkl}})^{\frac12}}+$$
$$+\sum_{j=1}^N\sum_{v\in R}c_{j,v}{\rm Expr}_n(f_v^{(1)})\cdot (\frac{T_{e_j}}{(1+\Delta_{{\rm Dunkl}})^{\frac12}}O_v+O_v\frac{T_{e_j}}{(1+\Delta_{{\rm Dunkl}})^{\frac12}})+$$
$$+{\rm Expr}_n(f^{(2)})\cdot (1+\Delta_{{\rm Dunkl}})^{-\frac12} +\sum_{v\in R}{\rm Expr}_n(f_v^{(2)})\cdot (1+\Delta_{{\rm Dunkl}})^{-\frac12}O_v+$$
$$+\sum_{v,v'\in R}{\rm Expr}_n(f_{v,v'}^{(2)})\cdot (1+\Delta_{{\rm Dunkl}})^{-\frac12}O_{v'}O_v.$$
In each summand, the first factor is bounded by inductive assumption and the second factor is obviously bounded. This yields the step of induction and, hence, completes the proof.
\end{proof}

\begin{proof}[Proof of Theorem \ref{smoothness theorem} \eqref{smootha}] The assertion follows from Lemma \ref{step of induction} and Hadamard 3 lines theorem.
\end{proof}

\begin{lemma}\label{second power lemma} For every $f\in C^{\infty}_c(\mathbb{R}^N),$ we have
$$M_f(1+\Delta_{{\rm Dunkl}})^{-1}\in \mathcal{L}_{\frac{N}{2},\infty}(L_2(\mathbb{R}^N,w(x)dx)).$$
\end{lemma}
\begin{proof} Set
$$\mathcal{D}=\sum_{j=1}^N\gamma_j\otimes T_{e_j},$$
where $\{\gamma_j\}_{j=1}^N$ are Pauli matrices. We have
$$1\otimes M_f(1+\Delta_{{\rm Dunkl}})^{-1}=(1\otimes M_f)(\mathcal{D}^2+1)^{-1}=(1\otimes M_f)(\mathcal{D}+i)^{-1}(\mathcal{D}-i)^{-1}=$$
$$=(\mathcal{D}+i)^{-1}(1\otimes M_f)(\mathcal{D}-i)^{-1}+[1\otimes M_f,(\mathcal{D}+i)^{-1}](\mathcal{D}-i)^{-1}=$$
$$=(\mathcal{D}+i)^{-1}(1\otimes M_f)(\mathcal{D}-i)^{-1}+(\mathcal{D}+i)^{-1}[\mathcal{D},1\otimes M_f](\mathcal{D}^2+1)^{-1}=$$
$$=(\mathcal{D}+i)^{-1}(1\otimes M_f)(\mathcal{D}-i)^{-1}+(\mathcal{D}+i)^{-1}(\sum_{j=1}^N\gamma_j\otimes [T_{e_j,M_f}])(\mathcal{D}^2+1)^{-1}=$$
$$=(\mathcal{D}+i)^{-1}(1\otimes M_f)(\mathcal{D}-i)^{-1}+\sum_{j=1}^N(\mathcal{D}+i)^{-1}(\gamma_j\otimes M_{\partial_jf})(\mathcal{D}^2+1)^{-1}+$$
$$+\sum_{j=1}^N\sum_{\alpha\in R_+}\alpha_j\kappa(\alpha)(\mathcal{D}+i)^{-1}(\gamma_j\otimes M_{f_{\alpha}}O_{\alpha})(\mathcal{D}^2+1)^{-1}.$$
Since $\Delta_{{\rm Dunkl}}$ commutes with $O_{\alpha},$ it follows that
$$1\otimes M_f(1+\Delta_{{\rm Dunkl}})^{-1}=$$
$$=(\mathcal{D}+i)^{-1}(1\otimes M_f)(\mathcal{D}-i)^{-1}+\sum_{j=1}^N(\mathcal{D}+i)^{-1}(1\otimes M_{\partial_jf})(\mathcal{D}^2+1)^{-1}\cdot(\gamma_j\otimes 1)+$$
$$+\sum_{j=1}^N\sum_{\alpha\in R_+}\alpha_j\kappa(\alpha)(\mathcal{D}+i)^{-1}(1\otimes M_{f_{\alpha}})(\mathcal{D}^2+1)^{-1}\cdot (\gamma_j\otimes O_{\alpha}).$$
Clearly,
$$\|(\mathcal{D}+i)^{-1}(1\otimes M_f)(\mathcal{D}-i)^{-1}\|_{\mathcal{L}_{\frac{N}{2},\infty}(\mathbb{C}^{2^{\lfloor\frac{N}{2}\rfloor}}\otimes L_2(\mathbb{R}^N,w(x)dx))}=$$
$$=\|(\Delta_{{\rm Dunkl}}+1)^{-\frac12}M_f(\Delta_{{\rm Dunkl}}+1)^{-\frac12}\|_{\mathcal{L}_{\frac{N}{2},\infty}(L_2(\mathbb{R}^N,w(x)dx))}\leq$$
$$=\|\Delta_{{\rm Dunkl}}^{-\frac12}M_f\Delta_{{\rm Dunkl}}^{-\frac12}\|_{\mathcal{L}_{\frac{N}{2},\infty}(L_2(\mathbb{R}^N,w(x)dx))}\leq$$
$$\leq 2^{\frac2N}\|M_{|f|^{\frac12}}\Delta_{{\rm Dunkl}}^{-\frac12}\|_{\mathcal{L}_{\frac{N}{2},\infty}(L_2(\mathbb{R}^N,w(x)dx))}^2\leq$$
$$\stackrel{Th.\ref{dunkl cwikel theorem}}{\leq}c_{N,\kappa}\||f|^{\frac12}\|_{L_N(\mathbb{R}^N)}^2=c_{N,\kappa}\|f\|_{L_{\frac{N}{2}}(\mathbb{R}^N)}$$
and, similarly,
$$\|(\mathcal{D}+i)^{-1}(1\otimes M_{\partial_jf})(\mathcal{D}^2+1)^{-1}\|_{\mathcal{L}_{\frac{N}{2},\infty}(\mathbb{C}^{2^{\lfloor\frac{N}{2}\rfloor}}\otimes L_2(\mathbb{R}^N,w(x)dx))}\leq c_{N,\kappa}\|\partial_jf\|_{L_{\frac{N}{2}}(\mathbb{R}^N)},$$
$$\|(\mathcal{D}+i)^{-1}(1\otimes M_{f_{\alpha}})(\mathcal{D}^2+1)^{-1}\|_{\mathcal{L}_{\frac{N}{2},\infty}(\mathbb{C}^{2^{\lfloor\frac{N}{2}\rfloor}}\otimes L_2(\mathbb{R}^N,w(x)dx))}\leq c_{N,\kappa}\|f_{\alpha}\|_{L_{\frac{N}{2}}(\mathbb{R}^N)}.$$
This completes the proof.
\end{proof}


\begin{lemma}\label{higher power lemma} For every $f\in C^{\infty}_c(\mathbb{R}^N),$ we have
$$M_f(1+\Delta_{{\rm Dunkl}})^{-\frac{n}{2}}\in \mathcal{L}_{\frac{N}{n},\infty}(L_2(\mathbb{R}^N,w(x)dx)).$$
\end{lemma}
\begin{proof} We prove the assertion by induction on $f.$ The $n$-th statement is "for every $f\in C^{\infty}_c(\mathbb{R}^N),$ we have
$$M_f(1+\Delta_{{\rm Dunkl}})^{-\frac{n}{2}}\in \mathcal{L}_{\frac{N}{n},\infty}(L_2(\mathbb{R}^N,w(x)dx))."$$
Base of induction (i.e. cases $n=1$ and $n=2$) are given by Theorem \ref{dunkl cwikel theorem} and Lemma \ref{second power lemma}. It suffices to prove the step of induction. 

We have
$$M_f(1+\Delta_{{\rm Dunkl}})^{-\frac{n+1}{2}}=$$
$$=(1+\Delta_{{\rm Dunkl}})^{-1}M_f(1+\Delta_{{\rm Dunkl}})^{-\frac{n-1}{2}}+[M_f,(1+\Delta_{{\rm Dunkl}})^{-1}](1+\Delta_{{\rm Dunkl}})^{-\frac{n-1}{2}}=$$
$$=(1+\Delta_{{\rm Dunkl}})^{-1}M_f(1+\Delta_{{\rm Dunkl}})^{-\frac{n-1}{2}}+(1+\Delta_{{\rm Dunkl}})^{-1}[\Delta_{{\rm Dunkl}},M_f](1+\Delta_{{\rm Dunkl}})^{-\frac{n+1}{2}}.$$
Choose $\phi\in C^{\infty}_c(\mathbb{R}^N)$ such that $(f\circ U_g)\cdot\phi=f\circ U_g$ for every $g\in G.$ It follows that
$$[\Delta_{{\rm Dunkl}},M_f]=[\Delta_{{\rm Dunkl}},M_f]M_{\phi}.$$
Thus,
$$M_f(1+\Delta_{{\rm Dunkl}})^{-\frac{n+1}{2}}=(1+\Delta_{{\rm Dunkl}})^{-1}M_f\cdot M_{\phi}(1+\Delta_{{\rm Dunkl}})^{-\frac{n-1}{2}}+$$
$$+(1+\Delta_{{\rm Dunkl}})^{-1}[\Delta_{{\rm Dunkl}},M_f]\cdot M_{\phi}(1+\Delta_{{\rm Dunkl}})^{-\frac{n}{2}}\cdot (1+\Delta_{{\rm Dunkl}})^{-\frac12}.$$
In the first summand on the right hand side, the first factor belongs to the ideal $\mathcal{L}_{\frac{N}{2},\infty}(L_2(\mathbb{R}^N,w(x)dx))$ by Lemma \ref{second power lemma} and the second one belongs to the ideal $\mathcal{L}_{\frac{N}{n-1},\infty}(L_2(\mathbb{R}^N,w(x)dx))$ by inductive assumption. In the second summand on the right hand side, the first factor belongs to the ideal $\mathcal{L}_{N,\infty}(L_2(\mathbb{R}^N,w(x)dx))$ by Lemma \ref{base of induction}, the second factor belongs to the ideal $\mathcal{L}_{\frac{N}{n},\infty}(L_2(\mathbb{R}^N,w(x)dx))$ by inductive assumption and the third one is obviously bounded. 
\end{proof}


\begin{lemma}\label{smoothness integer lemma} Let $k,l\in\mathbb{Z}_+$ with $l>k.$ For every $f\in C^{\infty}_c(\mathbb{R}^N),$ we have
$$T_{e_{m_1}}\cdots T_{e_{m_k}}M_f(1+\Delta_{{\rm Dunkl}})^{-\frac{l}{2}}$$
belongs to $\mathcal{L}_{\frac{N}{l-k},\infty}(L_2(\mathbb{R}^N,w(x)dx)).$
\end{lemma}
\begin{proof} We prove the assertion by induction on $k.$ Base of induction (i.e., the case $k=0$) is established in Lemma \ref{higher power lemma}.

We have
$$T_{e_{m_1}}\cdots T_{e_{m_k}}M_f(1+\Delta_{{\rm Dunkl}})^{-\frac{l}{2}}=T_{e_{m_1}}\cdots T_{e_{m_{k-1}}}M_fT_{e_{m_k}}(1+\Delta_{{\rm Dunkl}})^{-\frac{l}{2}}+$$
$$+T_{e_{m_1}}\cdots T_{e_{m_{k-1}}}[T_{e_{m_k}},M_f](1+\Delta_{{\rm Dunkl}})^{-\frac{l}{2}}=$$
$$=T_{e_{m_1}}\cdots T_{e_{m_{k-1}}}M_f(1+\Delta_{{\rm Dunkl}})^{-\frac{l-1}{2}}\times\frac{T_{e_{m_k}}}{(1+\Delta_{{\rm Dunkl}})^{\frac12}}+$$
$$+T_{e_{m_1}}\cdots T_{e_{m_{k-1}}}M_{\partial_{m_k}f}(1+\Delta_{{\rm Dunkl}})^{-\frac{l-1}{2}}\times (1+\Delta_{{\rm Dunkl}})^{-\frac12}+$$
$$+\frac12\sum_{v\in R}v_{m_k}\kappa(v)T_{e_{m_1}}\cdots T_{e_{m_{k-1}}}M_{\delta_vf}(1+\Delta_{{\rm Dunkl}})^{-\frac{l-1}{2}}\times  (1+\Delta_{{\rm Dunkl}})^{-\frac12}O_v.$$
In each summand on the right hand side, the first factor belongs to the ideal $\mathcal{L}_{\frac{N}{l-k},\infty}(L_2(\mathbb{R}^N,w(x)dx))$ by inductive assumption and the second factor is obviously bounded. This yields the step of induction and, hence, completes the proof.
\end{proof}


\begin{lemma}\label{smoothness final lemma} Let $s_1,s_2\in\mathbb{R}$ be such that $s_1+s_2>0.$ We have
$$(1+\Delta_{{\rm Dunkl}})^{-\frac{s_1}{2}} M_f(1+\Delta_{{\rm Dunkl}})^{-\frac{s_2}{2}}$$
belongs to $\mathcal{L}_{\frac{N}{s_1+s_2},\infty}(L_2(\mathbb{R}^N,w(x)dx)).$
\end{lemma}
\begin{proof} {\bf Step 1:} Let us prove the assertion for $s_1,s_2\in\mathbb{Z}.$

Suppose first $s_1\leq 0.$ Denote $k=-s_1\geq0$ and $l=s_2>k.$ Denote for brevity
$$V=(1+\Delta_{{\rm Dunkl}})^{\frac{k}{2}}(1+\sum_{j_1,\cdots,j_k=1}^N|T_{e_{j_1}}\cdots T_{e_{j_k}}|)^{-1}.$$
Obviously, $V$ is bounded. We have
$$(1+\Delta_{{\rm Dunkl}})^{\frac{k}{2}}=V\times (1+\sum_{m_1,\cdots,m_k=1}^N|T_{e_{m_1}}\cdots T_{e_{m_k}}|)=V+$$
$$+\sum_{m_1,\cdots,m_k=1}^NV\times{\rm sgn}(T_{e_{m_1}}\cdots T_{e_{m_k}})^{\ast}\times T_{e_{m_1}}\cdots T_{e_{m_k}}.$$
It, therefore, suffices to prove that operators
$$M_f(1+\Delta_{{\rm Dunkl}})^{-\frac{l}{2}},\quad T_{e_{m_1}}\cdots T_{e_{m_k}}M_f(1+\Delta_{{\rm Dunkl}})^{-\frac{l}{2}}$$
belong to $\mathcal{L}_{\frac{N}{s_1+s_2},\infty}(L_2(\mathbb{R}^N,w(x)dx)).$
This is established in Lemma \ref{higher power lemma} and Lemma \ref{smoothness integer lemma}, respectively. Hence, the claim is proved for $s_1,s_2\in\mathbb{Z}$ and $s_1\leq0.$

If $s_2\leq 0,$ then the claim follows from that for $s_2\leq0$ by taking the adjoint. If $s_1,s_2\geq0,$ then, according to Lemma \ref{smoothness integer lemma} the operators
$$(1+\Delta_{{\rm Dunkl}})^{-\frac{s_1+s_2}{2}} M_f,\quad M_f(1+\Delta_{{\rm Dunkl}})^{-\frac{s_1+s_2}{2}}$$
belong to $\mathcal{L}_{\frac{N}{s_1+s_2},\infty}(L_2(\mathbb{R}^N,w(x)dx)).$ The claim for $s_1,s_2\geq0$ follows now from the Hadamard 3 lines theorem.

{\bf Step 2:} The general case follows now from Step 1 by complex interpolation.
\end{proof}

\begin{proof}[Proof of Theorem \ref{smoothness theorem} \eqref{smoothb}] Choose $\phi\in C^{\infty}_c(\mathbb{R}^N)$ such that $(f\circ U_g)\cdot\phi=f\circ U_g$ for every $g\in G.$ It follows that
$$[\Delta_{{\rm Dunkl}},M_f]=[\Delta_{{\rm Dunkl}},M_f]M_{\phi}.$$
We have
$$(1+\Delta_{{\rm Dunkl}})^{-\frac{s_1}{2}}[\Delta_{{\rm Dunkl}},[\Delta_{{\rm Dunkl}},[\cdots,M_f]\cdots](1+\Delta_{{\rm Dunkl}})^{-\frac{s}{2}}=$$
$$=(1+\Delta_{{\rm Dunkl}})^{-\frac{s_1}{2}}[\Delta_{{\rm Dunkl}},[\Delta_{{\rm Dunkl}},[\cdots,M_f]\cdots]M_{\phi}(1+\Delta_{{\rm Dunkl}})^{-\frac{s}{2}}=$$
$$=(1+\Delta_{{\rm Dunkl}})^{-\frac{s_1}{2}}[\Delta_{{\rm Dunkl}},[\Delta_{{\rm Dunkl}},[\cdots,M_f]\cdots](1+\Delta_{{\rm Dunkl}})^{-\frac{n-s_1}{2}}\times$$
$$\times (1+\Delta_{{\rm Dunkl}})^{-\frac{s_1-n}{2}}M_{\phi}(1+\Delta_{{\rm Dunkl}})^{-\frac{s_2}{2}}.$$
The first factor is bounded by Theorem \ref{smoothness theorem} \eqref{smootha}. The second factor belongs to $\mathcal{L}_{\frac{N}{s_1+s_2-n},\infty}(L_2(\mathbb{R}^N,w(x)dx))$ by Lemma \ref{smoothness final lemma}.
\end{proof}

\section{Approximation theorem}

\begin{theorem}\label{final approximation theorem} For every $f\in C^{\infty}_c(\mathbb{R}^N),$ we have
$$M_{w^{\frac12}}[R_{{\rm Dunkl},j},M_f]M_{w^{-\frac12}}\in V\Delta^{-\frac12}+(\mathcal{L}_{N,\infty})_0(L_2(\mathbb{R}^N)),$$
$$V=M_{\partial_jf}+\frac12\sum_{v\in R}v_k\kappa(v)M_{\delta_vf}O_v-\sum_{k=1}^NR_jM_{\partial_kf}R_k-$$
$$-\frac14R_j\cdot\sum_{k=1}^N\sum_{v\in R}v_k\kappa(v)M_{\delta_vf}(R_kO_v+O_vR_k).$$
\end{theorem}


\subsection{First approximation theorem}

\begin{theorem}\label{first approximation theorem} For every $f\in C^{\infty}_c(\mathbb{R}^N),$ we have
$$[(1+\Delta_{{\rm Dunkl}})^{\frac12},M_f](1+\Delta_{{\rm Dunkl}})^{-\frac12}-\frac12[\Delta_{{\rm Dunkl}},M_f](1+\Delta_{{\rm Dunkl}})^{-1}\in$$
$$\in\mathcal{L}_{\frac{2N}{3},\infty}(L_2(\mathbb{R}^N,w(x)dx)).$$
\end{theorem}

\begin{lemma}\label{cgrs-like abstract lemma} Let $B\geq 1.$ Suppose
$$B^{-1}[B^2,[B^2,A]]B^{-\frac{5}{2}}\in\mathcal{L}_{\frac{2N}{3},\infty}.$$
We have
$$[B,A]B^{-1}-\frac12[B^2,A]B^{-2}\in\mathcal{L}_{\frac{2N}{3},\infty}.$$
\end{lemma}
\begin{proof} We have
$$B=\frac{2}{\pi}\int_0^{\infty}\frac{B^2d\lambda}{\lambda^2+B^2}.$$
Hence,
$$[B,A]B^{-1}=\frac{2}{\pi}\int_0^{\infty}[\frac{B^2}{\lambda^2+B^2},A]B^{-1}d\lambda.$$
Obviously,
$$[\frac{B^2}{\lambda^2+B^2},A]B^{-1}=-[\frac{\lambda^2}{\lambda^2+B^2},A]B^{-1}=\frac{\lambda^2}{\lambda^2+B^2}[B^2,A]\frac1{B(\lambda^2+B^2)}=$$
$$=[B^2,A]\frac{\lambda^2}{B(\lambda^2+B^2)^2}-\frac{\lambda^2}{\lambda^2+B^2}[B^2,[B^2,A]]\frac1{B(\lambda^2+B^2)^2}.$$
Obviously,
$$\int_0^{\infty}[B^2,A]\frac{\lambda^2}{B(\lambda^2+B^2)^2}d\lambda=[B^2,A]B^{-1}\cdot \int_0^{\infty}\frac{\lambda^2}{(\lambda^2+B^2)^2}d\lambda=$$
$$=[B^2,A]B^{-2}\cdot \int_0^{\infty}\frac{\lambda^2}{(\lambda^2+1)^2}d\lambda=\frac{\pi}{4}[B^2,A]B^{-2}.$$

On the other hand,
$$\|\frac{\lambda^2}{\lambda^2+B^2}[B^2,[B^2,A]]\frac1{B(\lambda^2+B^2)^2}\|_{\frac{2N}{3},\infty}\leq$$
$$\leq\|\frac{\lambda^2B}{\lambda^2+B^2}\|_{\infty}\|B^{-1}[B^2,[B^2,A]]B^{-\frac{5}{2}}\|_{\frac{2N}{3},\infty}\|\frac{B^{\frac32}}{(\lambda^2+B^2)^2}\|_{\infty}.$$
Next,
$$\|\frac{\lambda^2B}{\lambda^2+B^2}\|_{\infty}\leq\sup_{t\geq1}\frac{\lambda^2t}{\lambda^2+t^2}\leq\lambda$$
and
$$\|\frac{B^{\frac32}}{(\lambda^2+B^2)^2}\|_{\infty}\leq\sup_{t\geq1}\frac{t^{\frac32}}{(\lambda^2+t^2)^2}\leq(\lambda^2+1)^{-\frac54}.$$
Thus,
$$\|\frac{\lambda^2}{\lambda^2+B^2}[B^2,[B^2,A]]\frac1{B(\lambda^2+B^2)^2}\|_{\frac{2N}{3},\infty}\leq$$
$$\leq\frac{\lambda}{(\lambda^2+1)^{\frac54}}\|B^{-1}[B^2,[B^2,A]]B^{-\frac{5}{2}}\|_{\frac{2N}{3},\infty}.$$
Hence, the integral
$$\int_0^{\infty}\frac{\lambda^2}{\lambda^2+B^2}[B^2,[B^2,A]]\frac1{B(\lambda^2+B^2)^2}d\lambda$$
converges absolutely in $\mathcal{L}_{\frac{2N}{3},\infty}$ and the assertion follows.
\end{proof}

\begin{proof}[Proof of Theorem \ref{first approximation theorem}] The assertion follows from Theorem \ref{smoothness theorem} and Lemma \ref{cgrs-like abstract lemma}.
\end{proof}

\subsection{Second approximation theorem}

\begin{theorem}\label{second approximation theorem} If $f\in C^{\infty}_c(\mathbb{R}^N),$ then
$$M_f(1+\Delta_{{\rm Dunkl}})^{-\frac12}-M_f(1+M_{w^{-\frac12}}\Delta M_{w^{\frac12}})^{-\frac12}\in (\mathcal{L}_{N,\infty})_0(L_2(\mathbb{R}^N,w(x)dx)).$$
\end{theorem}

\begin{lemma}\label{s10 abstract lemma} Let $A,B_1,B_2$ be such that $B_1,B_2\geq1.$ Suppose that
\begin{enumerate}[{\rm (i)}]
\item\label{s10aa} $A(B_2-B_1)B_2^{-\frac12},$ $[B_1,A]B_1^{-\frac12}$ and $[B_1,A]B_2^{-\frac12}$ are bounded;
\item\label{s10ab} there exists $A'$ with $B_1^{-\frac34}A'\in\mathcal{L}_{\frac{2N}{3},\infty}$ such that $A=A'A$ and $[B_1,A]=A'[B_1,A];$
\end{enumerate}
We have
$$AB_1^{-\frac12}-AB_2^{-\frac12}\in\mathcal{L}_{\frac{2N}{3},\infty}.$$
\end{lemma}
\begin{proof} We have
$$B_k^{-\frac12}=\frac{2}{\pi}\int_0^{\infty}\frac{d\lambda}{\lambda^2+B_k},\quad k=1,2.$$
Thus,
$$AB_1^{-\frac12}-AB_2^{-\frac12}=\int_0^{\infty}A(\frac1{\lambda^2+B_1}-\frac1{\lambda^2+B_2})d\lambda.$$
Obviously,
$$A(\frac1{\lambda^2+B_1}-\frac1{\lambda^2+B_2})=A\frac1{\lambda^2+B_1}(B_2-B_1)\frac1{\lambda^2+B_2}=$$
$$=\frac1{\lambda^2+B_1}A(B_2-B_1)\frac1{\lambda^2+B_2}+\frac1{\lambda^2+B_1}[B_1,A]\frac1{\lambda^2+B_1}(B_2-B_1)\frac1{\lambda^2+B_2}=$$
$$=\frac1{\lambda^2+B_1}A(B_2-B_1)\frac1{\lambda^2+B_2}+\frac1{\lambda^2+B_1}[B_1,A]\frac1{\lambda^2+B_1}-\frac1{\lambda^2+B_1}[B_1,A]\frac1{\lambda^2+B_2}.$$

Consider, for example, the first summand. We write
$$\frac1{\lambda^2+B_1}A(B_2-B_1)\frac1{\lambda^2+B_2}=\frac1{\lambda^2+B_1}A'\cdot A(B_2-B_1)\frac1{\lambda^2+B_2}=$$
$$=\frac{B_1^{\frac34}}{\lambda^2+B_1}\cdot B_1^{-\frac34}A'\cdot A(B_2-B_1)B_2^{-\frac12}\cdot\frac{B_2^{\frac12}}{\lambda^2+B_2}.$$
Thus,
$$\|\frac1{\lambda^2+B_1}A(B_2-B_1)\frac1{\lambda^2+B_2}\|_{\frac{3N}{4},\infty}\leq$$
$$\leq\|\frac{B_1^{\frac34}}{\lambda^2+B_1}\|_{\infty}\cdot\|B_1^{-\frac34}A'\|_{\frac{3N}{4},\infty}\cdot\|A(B_2-B_1)B_2^{-\frac12}\|_{\infty}\cdot\|\frac{B_2^{\frac12}}{\lambda^2+B_2}\|_{\infty}\leq$$
$$\leq\sup_{t\geq1}\frac{t^{\frac34}}{\lambda^2+t}\cdot\|B_1^{-\frac34}A'\|_{\frac{3N}{4},\infty}\cdot\|A(B_2-B_1)B_2^{-\frac12}\|_{\infty}\cdot\sup_{t\geq1}\frac{t^{\frac12}}{\lambda^2+t}\leq$$
$$\leq(\lambda^2+1)^{-\frac14}\cdot\|B_1^{-\frac34}A'\|_{\frac{3N}{4},\infty}\cdot\|A(B_2-B_1)B_2^{-\frac12}\|_{\infty}\cdot(\lambda^2+1)^{-\frac12}.$$
Hence, the first summand is absolutely integrable in $\mathcal{L}_{\frac{3N}{4},\infty}.$ Similarly, so are the other summands. This completes the proof.
\end{proof}

\begin{lemma}\label{s10 verification lemma} Let $P$ be a differential operator of degree $1$ with smooth compactly supported coefficients. Suppose that the coefficients are supported away from the hyperplanes $\{\{x:\ \langle x,v\rangle=0\}\}_{v\in R}.$ The operators
$$P(1+\Delta_{{\rm Dunkl}})^{-\frac12},\quad P(1+M_{w^{-\frac12}}\Delta M_{w^{\frac12}})^{-\frac12}$$
are bounded on $L_2(\mathbb{R}^N,w(x)dx).$
\end{lemma}
\begin{proof} By assumption, $P=M_{w^{-\frac12}}P'M_{w^{\frac12}},$ where $P'$ is a differential operator of degree $1$ with smooth compactly supported coefficients. Thus,
$$P(1+M_{w^{-\frac12}}\Delta M_{w^{\frac12}})^{-\frac12}=M_{w^{-\frac12}}\cdot P'(1+\Delta)^{-\frac12}\cdot M_{w^{\frac12}}.$$
Boundedness of the right hand side is obvious.

Similarly,
$$P=\sum_{j=1}^NM_{f_j}T_{e_j}+\sum_{g\in G}M_{h_g}U_g,$$
where functions $\{f_j\}_{j=1}^N$ and $\{h_g\}_{g\in G}$ are smooth and compactly supported. Since the operators $\{U_g\}_{g\in G}$ commute with $\Delta_{{\rm Dunkl}},$ it follows that
$$P(1+\Delta_{{\rm Dunkl}})^{-\frac12}=\sum_{j=1}^NM_{f_j}\frac{T_{e_j}}{(1+\Delta_{{\rm Dunkl}})^{\frac12}}+\sum_{g\in G}M_{h_g}\frac{T_{e_j}}{(1+\Delta_{{\rm Dunkl}})^{\frac12}}U_g.$$
Boundedness of the right hand side is obvious.
\end{proof}

\begin{proof}[Proof of Theorem \ref{second approximation theorem}]

{\bf Step 1:} Suppose first that $f$ is supported away from the hyperplanes $\{\{x:\ \langle x,v\rangle=0\}\}_{v\in R}.$

We apply Lemma \ref{s10 abstract lemma} with $A=M_f,$ $B_1=1+M_{w^{-\frac12}}\Delta M_{w^{\frac12}}$ and $B_2=1+\Delta_{{\rm Dunkl}}.$

Choose $\phi\in C^{\infty}_c(\mathbb{R}^N)$ such that $f=f\phi$ and set $A'=M_{\phi}.$ We have
$$M_{\phi}[M_{w^{-\frac12}}\Delta M_{w^{\frac12}},M_f]=[M_{w^{-\frac12}}\Delta M_{w^{\frac12}},M_f].$$
Thus, $A=A'A$ and $[B_1,A]=A'[B_1,A].$ It follows from the usual Cwikel estimates that $B_1^{-\frac34}A'\in\mathcal{L}_{\frac{2N}{3},\infty}(L_2(\mathbb{R}^N,w(x)dx)).$ This verifies the assumption \eqref{s10ab} in Lemma \ref{s10 abstract lemma}.

Note that $[B_1,A]$ is a differential operator $P$ of degree $1$ with smooth compactly supported coefficients. Furthermore, the coefficients are supported away from the hyperplanes $\{\{x:\ \langle x,v\rangle=0\}\}_{v\in R}.$ By Lemma \ref{s10 verification lemma}, the operators
$$[B_1,A]B_1^{-\frac12}=P(1+M_{w^{-\frac12}}\Delta M_{w^{\frac12}})^{-\frac12},\quad [B_1,A]B_2^{-\frac12}=P(1+\Delta_{{\rm Dunkl}})^{-\frac12},$$
are bounded on $L_2(\mathbb{R}^N,w(x)dx).$ Note that there exist differential operators $\{P_g\}_{g\in G}$ of degree $1$ with smooth compactly supported coefficients such that the coefficients are supported away from the hyperplanes $\{\{x:\ \langle x,v\rangle=0\}\}_{v\in R}$ and such that
$$A(B_2-B_1)=\sum_{g\in G}P_gU_g.$$
Since $\{U_g\}_{g\in G}$ commute with $\Delta_{{\rm Dunkl}},$ it follows that
$$A(B_2-B_1)=\sum_{g\in G}P_g(1+\Delta_{{\rm Dunkl}})^{-\frac12}\cdot U_g$$
is bounded on $L_2(\mathbb{R}^N,w(x)dx).$ This verifies the assumption \eqref{s10aa} in Lemma \ref{s10 abstract lemma}.

Using Lemma \ref{s10 abstract lemma}, we infer the claim of Step 1.

{\bf Step 2:} Choose a sequence $\{\phi_n\}_{n\geq1}\subset C^{\infty}_c(\mathbb{R}^N)$ such that (a) $0\leq\phi_n\leq1$ for $n\geq1$ (b) $\phi_n$ is supported away from the hyperplanes $\{x:\ \langle x,v\rangle=0\},$ $v\in R,$ for $n\geq 1$ (c) ${\rm supp}(\phi_n)\uparrow\mathbb{R}^N$ as $n\to\infty.$

By Step 1,
$$M_{f\phi_n}(1+\Delta_{{\rm Dunkl}})^{-\frac12}-M_{f\phi_n}(1+M_{w^{-\frac12}}\Delta M_{w^{\frac12}})^{-\frac12}\in (\mathcal{L}_{N,\infty})_0(L_2(\mathbb{R}^N,w(x)dx))$$
for every $n\geq1.$

By Theorem \ref{dunkl cwikel theorem},
$$M_{f\phi_n}(1+\Delta_{{\rm Dunkl}})^{-\frac12}\to M_f(1+\Delta_{{\rm Dunkl}})^{-\frac12}$$
in $\mathcal{L}_{N,\infty}(L_2(\mathbb{R}^N,w(x)dx))$ as $n\to\infty.$ By the usual Cwikel estimate,
$$M_{f\phi_n}(1+M_{w^{-\frac12}}\Delta M_{w^{\frac12}})^{-\frac12}\to M_f(1+M_{w^{-\frac12}}\Delta M_{w^{\frac12}})^{-\frac12}$$
in $\mathcal{L}_{N,\infty}(L_2(\mathbb{R}^N,w(x)dx))$ as $n\to\infty.$ Thus
$$\lim_{n\to\infty}M_{f\phi_n}(1+\Delta_{{\rm Dunkl}})^{-\frac12}-M_{f\phi_n}(1+M_{w^{-\frac12}}\Delta M_{w^{\frac12}})^{-\frac12}=$$
$$=M_f(1+\Delta_{{\rm Dunkl}})^{-\frac12}-M_f(1+M_{w^{-\frac12}}\Delta M_{w^{\frac12}})^{-\frac12},$$
where the limit is understood in $\mathcal{L}_{N,\infty}(L_2(\mathbb{R}^N,w(x)dx)).$ Since the separable part $(\mathcal{L}_{N,\infty})_0(L_2(\mathbb{R}^N,w(x)dx))$ is closed in $\mathcal{L}_{N,\infty}(L_2(\mathbb{R}^N,w(x)dx)),$ the assertion follows.
\end{proof}




\subsection{Proof of Theorem \ref{final approximation theorem}}



\begin{lemma}\label{trivial approximation lemma} We have
$$[R_{{\rm Dunkl},j},M_f]-[\frac{T_{e_j}}{(1+\Delta_{{\rm Dunkl}})^{\frac12}},M_f]\in \mathcal{L}_{\frac{N}{2},\infty}(L_2(\mathbb{R}^N,w(x)dx)).$$
\end{lemma}
\begin{proof} Set 
$$h(x)=1-\frac{|x|}{(1+|x|^2)^{\frac12}},\quad x\in\mathbb{R}^N.$$
We, obviously, have
$$R_{{\rm Dunkl},j}-\frac{T_{e_j}}{(1+\Delta_{{\rm Dunkl}})^{\frac12}}=R_{{\rm Dunkl},j}\cdot h(\nabla_{{\rm Dunkl}}).$$
Thus,
$$[R_{{\rm Dunkl},j},M_f]-[\frac{T_{e_j}}{(1+\Delta_{{\rm Dunkl}})^{\frac12}},M_f]=$$
$$=R_{{\rm Dunkl},j}\cdot h(\nabla_{{\rm Dunkl}})M_f-M_fh(\nabla_{{\rm Dunkl}})\cdot R_{{\rm Dunkl},j}=$$
$$=R_{{\rm Dunkl},j}\cdot h(\nabla_{{\rm Dunkl}})(1+\Delta_{{\rm Dunkl}})\cdot (1+\Delta_{{\rm Dunkl}})^{-1} M_f-$$
$$-M_f(1+\Delta_{{\rm Dunkl}})^{-1}\cdot (1+\Delta_{{\rm Dunkl}})h(\nabla_{{\rm Dunkl}})\cdot R_{{\rm Dunkl},j}.$$
In the first summand on the right hand side, the first and the second factors are bounded and the third one belongs to $\mathcal{L}_{\frac{N}{2},\infty}(L_2(\mathbb{R}^N,w(x)dx))$ by Lemma \ref{higher power lemma}. In the second summand on the right hand side, the second and the third factors are bounded and the first one belongs to $\mathcal{L}_{\frac{N}{2},\infty}(L_2(\mathbb{R}^N,w(x)dx))$ by Lemma \ref{higher power lemma}. 
\end{proof}

\begin{lemma}\label{s11 abstract lemma} Let $\{B_j\}_{j=1}^N$ and $\{D_j\}_{j=1}^N$ be commuting tuples of self-adjoint operators. Denote $\langle B\rangle=(1+\sum_{j=1}^NB_j^2)^{\frac12}$ and $\langle D\rangle=(1+\sum_{j=1}^ND_j^2)^{\frac12}.$ Suppose that there exists a sequence $\{\phi_n\}_{n\geq1}\subset C^{\infty}_c(\mathbb{R}^N)$ with
\begin{enumerate}[{\rm (i)}]
\item $0\leq\phi_n\leq1$ for $n\geq1$ and ${\rm supp}(\phi_n)\uparrow\mathbb{R}^N$ as $n\to\infty;$
\item 
$$M_{\phi_n}\langle B\rangle^{-1}-M_{\phi_n}\langle D\rangle^{-1}\in(\mathcal{L}_{N,\infty})_0(L_2(\mathbb{R}^N,w(x)dx)),\quad n\geq1;$$
\item $M_{\phi_n}(B_k-D_k)$ is bounded\footnote{By a constant which depends also on $n.$} and compactly supported for $n\geq 1$ and $1\leq k\leq K;$
\item $[D_k,M_{f\phi_n}]$ is bounded\footnote{By a constant which depends also on $n.$} and compactly supported for $n\geq 1$ and $1\leq k\leq K;$
\item $[\langle D\rangle^{-1},M_{f\phi_n}]\in(\mathcal{L}_{N,\infty})_0(\mathbb{R}^N,w(x)dx)$ for every $n\geq1;$
\item
$$\|M_h\langle B\rangle^{-1}\|_{\mathcal{L}_{N,\infty}(\mathbb{R}^N,w(x)dx)}\leq c_N\|h\|_{L_N(\mathbb{R}^N)},\quad h\in L_N(\mathbb{R}^N),$$
$$\|M_h\langle D\rangle^{-1}\|_{\mathcal{L}_{N,\infty}(\mathbb{R}^N,w(x)dx)}\leq c_N\|h\|_{L_N(\mathbb{R}^N)},\quad h\in L_N(\mathbb{R}^N);$$
\end{enumerate}
We have
$$\frac{B_j}{\langle B\rangle}M_f\frac{B_k}{\langle B\rangle^2}-\frac{D_j}{\langle D\rangle}M_f\frac{D_k}{\langle D\rangle^2}\in(\mathcal{L}_{N,\infty})_0(L_2(\mathbb{R}^N,w(x)dx)).$$
\end{lemma}
\begin{proof} Choose real-valued $\psi_n\in C^{\infty}_c(\mathbb{R}^N)$ such that
$$M_{\phi_n}(B_k-D_k)=M_{\phi_n}(B_k-D_k)\cdot M_{\psi_n},\quad 1\leq k\leq K,$$
$$[M_{f\phi_n},D_k]=M_{\psi_n}\cdot [M_{f\phi_n},D_k]\cdot M_{\psi_n},\quad 1\leq k\leq K.$$
We have
$$\frac{B_j}{\langle B\rangle}M_{f\phi_n}\frac{B_k}{\langle B\rangle^2}-\frac{D_j}{\langle D\rangle}M_{f\phi_n}\frac{D_k}{\langle D\rangle^2}=\frac{B_j}{\langle B\rangle}\cdot M_f\Big(M_{\phi_n}\langle B\rangle^{-1}-M_{\phi_n}\langle D\rangle^{-1}\Big)\cdot\frac{B_k}{\langle B\rangle}+$$
$$+\frac{B_j}{\langle B\rangle}\cdot [M_{f\phi_n},\langle D\rangle^{-1}]\cdot\frac{B_k}{\langle B\rangle}+\frac{B_j}{\langle B\rangle}\cdot\langle D\rangle^{-1}M_f\cdot M_{\phi_n}(B_k-D_k)\cdot M_{\psi_n}\langle B\rangle^{-1}+$$
$$+\frac{B_j}{\langle B\rangle}\cdot\langle D\rangle^{-1}M_{\psi_n}\cdot[M_{f\phi_n},D_k]\cdot M_{\psi_n}\langle B\rangle^{-1}+\frac{B_j}{\langle B\rangle}\frac{D_k}{\langle D\rangle}\cdot M_f\Big(M_{\phi_n}\langle B\rangle^{-1}-M_{\phi_n}\langle D\rangle^{-1}\Big)+$$
$$+\frac{B_j}{\langle B\rangle}\cdot\langle D\rangle^{-1}M_{\psi_n}\cdot[D_k,M_{f\phi_n}]\cdot M_{\psi_n}\langle D\rangle^{-1}+\frac{B_j}{\langle B\rangle}\cdot[\langle D\rangle^{-1},M_{f\phi_n}]\cdot\frac{D_k}{\langle D\rangle}+$$
$$+\langle B\rangle^{-1}M_{\psi_n}\cdot (B_j-D_j)M_{\phi_n}\cdot M_f\frac{D_k}{\langle D\rangle^2}+$$
$$+\langle B\rangle^{-1}M_{\psi_n}\cdot [D_j,M_{f\phi_n}]\cdot M_{\psi_n}\frac{D_k}{\langle D\rangle^2}+(\langle B\rangle^{-1}-\langle D\rangle^{-1})M_{\phi_n}\cdot M_f\frac{D_jD_k}{\langle D\rangle^2}+$$
$$+\langle D\rangle^{-1}M_{\psi_n}\cdot[M_{f\phi_n},D_j]\cdot M_{\psi_n}\frac{D_k}{\langle D\rangle^2}.$$
It follows now from the assumptions and from H\"older inequality that
$$\frac{B_j}{\langle B\rangle}M_{f\phi_n}\frac{B_k}{\langle B\rangle^2}-\frac{D_j}{\langle D\rangle}M_{f\phi_n}\frac{D_k}{\langle D\rangle^2}\in(\mathcal{L}_{N,\infty})_0(L_2(\mathbb{R}^N,w(x)dx)).$$
Obviously,
$$\Big\|\frac{B_j}{\langle B\rangle}M_{f\phi_n}\frac{B_k}{\langle B\rangle^2}-\frac{B_j}{\langle B\rangle}M_f\frac{B_k}{\langle B\rangle^2}\Big\|_{\mathcal{L}_{N,\infty}(L_2(\mathbb{R}^N,w(x)dx))}\leq$$
$$\leq\Big\|M_{f(1-\phi_n)}\langle B\rangle^{-1}\Big\|_{\mathcal{L}_{N,\infty}(L_2(\mathbb{R}^N,w(x)dx))}\leq c_N\|f(1-\phi_n)\|_{L_N(\mathbb{R}^N)},$$
$$\Big\|\frac{D_j}{\langle D\rangle}M_{f\phi_n}\frac{D_k}{\langle D\rangle^2}-\frac{D_j}{\langle D\rangle}M_f\frac{D_k}{\langle D\rangle^2}\Big\|_{\mathcal{L}_{N,\infty}(L_2(\mathbb{R}^N,w(x)dx))}\leq$$
$$\leq\Big\|M_{f(1-\phi_n)}\langle D\rangle^{-1}\Big\|_{\mathcal{L}_{N,\infty}(L_2(\mathbb{R}^N,w(x)dx))}\leq c_N\|f(1-\phi_n)\|_{L_N(\mathbb{R}^N)}.$$
Thus,
$$\lim_{n\to\infty}\frac{B_j}{\langle B\rangle}M_{f\phi_n}\frac{B_k}{\langle B\rangle^2}-\frac{D_j}{\langle D\rangle}M_{f\phi_n}\frac{D_k}{\langle D\rangle^2}=$$
$$=\frac{B_j}{\langle B\rangle}M_f\frac{B_k}{\langle B\rangle^2}-\frac{D_j}{\langle D\rangle}M_f\frac{D_k}{\langle D\rangle^2},$$
where the limit is take in $\mathcal{L}_{N,\infty}(L_2(\mathbb{R}^N,w(x)dx)).$ Since $(\mathcal{L}_{N,\infty})_0(L_2(\mathbb{R}^N,w(x)dx))$ is closed in $\mathcal{L}_{N,\infty}(L_2(\mathbb{R}^N,w(x)dx)),$ the assertion follows.
\end{proof}

\begin{lemma}\label{s11 main lemma} For every $f\in C^{\infty}_c(\mathbb{R}^N),$ 
$$\frac{T_{e_j}}{(1+\Delta_{{\rm Dunkl}})^{\frac12}}M_f\frac{T_{e_k}}{1+\Delta_{{\rm Dunkl}}}O_v\in$$
$$\in M_{w^{-\frac12}}R_jM_fR_kO_v\Delta^{-\frac12}M_{w^{\frac12}}+(\mathcal{L}_{N,\infty})_0(L_2(\mathbb{R}^N,w(x)dx)),$$
$$\frac{T_{e_j}}{(1+\Delta_{{\rm Dunkl}})^{\frac12}}M_fO_v\frac{T_{e_j}}{1+\Delta_{{\rm Dunkl}}}\in $$
$$\in M_{w^{-\frac12}}R_jM_fO_vR_k\Delta^{-\frac12}M_{w^{\frac12}}+(\mathcal{L}_{N,\infty})_0(L_2(\mathbb{R}^N,w(x)dx)).$$
\end{lemma}
\begin{proof} The assertion follows by applying Lemma \ref{s11 abstract lemma} to the tuples $(T_{e_1},\cdots,T_{e_N})$ and $(M_{w^{-\frac12}}\partial_1M_{w^{\frac12}},\cdot,M_{w^{-\frac12}}\partial_NM_{w^{\frac12}})$ and sequence $\{\phi_n\}_{n\geq1}\subset C^{\infty}_c(\mathbb{R}^N)$ supported away from the hyperplanes $\{\{x:\ \langle x,v\rangle=0\}\}_{v\in R}.$
\end{proof}



\begin{proof}[Proof of Theorem \ref{final approximation theorem}] We have
$$[\frac{T_{e_j}}{(1+\Delta_{{\rm Dunkl}})^{\frac12}},M_f]=$$
$$=[T_{e_j},M_f](1+\Delta_{{\rm Dunkl}})^{-\frac12}-\frac{T_{e_j}}{(1+\Delta_{{\rm Dunkl}})^{\frac12}}\cdot [(1+\Delta_{{\rm Dunkl}})^{\frac12},M_f](1+\Delta_{{\rm Dunkl}})^{-\frac12}.$$
Using Lemma \ref{trivial approximation lemma} and Theorem \ref{first approximation theorem}, we write
$$[R_{{\rm Dunkl},j},M_f]\in [T_{e_j},M_f](1+\Delta_{{\rm Dunkl}})^{-\frac12}-$$
$$-\frac12\cdot\frac{T_{e_j}}{(1+\Delta_{{\rm Dunkl}})^{\frac12}}\cdot [\Delta_{{\rm Dunkl}},M_f](1+\Delta_{{\rm Dunkl}})^{-1}\in(\mathcal{L}_{N,\infty})_0(L_2(\mathbb{R}^N,w(x)dx)).$$
It follows from Lemma \ref{commutator with delta} and Lemma \ref{higher power lemma} that
$$[\Delta_{{\rm Dunkl}},M_f](1+\Delta_{{\rm Dunkl}})^{-1}-2\sum_{k=1}^NM_{\partial_kf}\frac{T_{e_k}}{1+\Delta_{{\rm Dunkl}}}-$$
$$-\frac12\sum_{k=1}^N\sum_{v\in R}v_k\kappa(v)M_{\delta_vf}\big(\frac{T_{e_k}}{1+\Delta_{{\rm Dunkl}}}O_v+O_v\frac{T_{e_k}}{1+\Delta_{{\rm Dunkl}}}\big)\in (\mathcal{L}_{N,\infty})_0(L_2(\mathbb{R}^N,w(x)dx)).$$
Thus,
$$[R_{{\rm Dunkl},j},M_f]\in M_{\partial_jf}(1+\Delta_{{\rm Dunkl}})^{-\frac12}+\frac12\sum_{v\in R}v_j\kappa(v)M_{\delta_vf}(1+\Delta_{{\rm Dunkl}})^{-\frac12}O_v-$$
$$-\frac{T_{e_j}}{(1+\Delta_{{\rm Dunkl}})^{\frac12}}\cdot \sum_{k=1}^NM_{\partial_kf}\frac{T_{e_k}}{1+\Delta_{{\rm Dunkl}}}-$$
$$-\frac14\frac{T_{e_j}}{(1+\Delta_{{\rm Dunkl}})^{\frac12}}\cdot\sum_{k=1}^N\sum_{v\in R}v_k\kappa(v)M_{\delta_vf}\big(\frac{T_{e_k}}{1+\Delta_{{\rm Dunkl}}}O_v+O_v\frac{T_{e_k}}{1+\Delta_{{\rm Dunkl}}}\big)\in$$
$$\in(\mathcal{L}_{N,\infty})_0(L_2(\mathbb{R}^N,w(x)dx)).$$
The assertion follows from Lemma \ref{s11 main lemma}.
\end{proof}



\section{Spectral asymptotic formula}
\setcounter{equation}{0}

We have a natural action $\alpha$ of $G$ on $B(L_2(\mathbb{R}^N))$ and, therefore, on $\Pi$ --- the $C^{\ast}$-algebra of operators admitting principal symbol. {\color{red} write precisely if it is $T\to U_g^{-1}TU_g$ or $T\to U_gTU_g^{-1}.$}

Let ${\rm sym}:\Pi\to(C_0+\mathbb{C})(\mathbb{R}^N)\otimes_{{\rm min}}C(\mathbb{S}^{N-1})$ be the principal symbol mapping.

Define the action $\beta$ of $G$ on $(C_0+\mathbb{C})(\mathbb{R}^N)\otimes_{{\rm min}}C(\mathbb{S}^{N-1})$ by setting
$$(\beta(g))(f\otimes h)=(f\circ g)\otimes (h\circ g).$$
{\color{red} check whether there should be $g$ or $g^{-1}$ in each of the tensor factors.}

We have equivariance relation {\color{red} check it!}
$${\rm sym}(\alpha(g)(T))=(\beta(g))({\rm sym}(T)).$$

Clearly, ${\rm sym}$ is surjective, ${\rm ker}({\rm sym})$ is $\alpha$-invariant. Hence, there exists a natural $\ast$-homomorphism $$\pi\rtimes G:\Pi\rtimes_{\alpha}G\to ((C_0+\mathbb{C})(\mathbb{R}^N)\otimes_{{\rm min}}C(\mathbb{S}^{N-1}))\rtimes_{\beta}G.$$

Next, the standard trace $\int_{\mathbb{R}^N}\otimes\int_{\mathbb{S}^{N-1}}$ on $(C_0+\mathbb{C})(\mathbb{R}^N)\otimes_{{\rm min}}C(\mathbb{S}^{N-1})$ is $\beta$-invariant. Hence, there exists a natural trace 
$$(\int_{\mathbb{R}^N}\otimes\int_{\mathbb{S}^{N-1}})\rtimes G$$
on $((C_0+\mathbb{C})(\mathbb{R}^N)\otimes_{{\rm min}}C(\mathbb{S}^{N-1}))\rtimes_{\beta}G.$


\begin{theorem} Let $N>2.$ Let $\Pi$ be the $C^{\ast}$-algebra of operators admitting principal symbol. For every $T\in\Pi\rtimes_{\alpha}G$ compactly supported from the right, we have
$$\lim_{t\to\infty}t^{\frac1N}\mu(t,T\Delta^{-\frac12})=c_{\alpha}\|({\rm sym}\rtimes G)(T)\|_{L_N(L_{\infty}(\mathbb{R}^N)\bar{\otimes}L_{\infty}(\mathbb{S}^{N-1}))\rtimes_{\beta})}.$$
\end{theorem}


\appendix

\section{Zhijie lemma}
\setcounter{equation}{0}

\begin{lemma}\label{lem:Cwikel-dyadic-distribution}
Let $1<s<\infty$, and let $\mathbb D^0$ be the standard dyadic grid
in $\mathbb R^N$. Suppose that
$
F\in L^s(\mathbb R^N).
$
Then there exists a constant $C>0$ such that
\[
\left\|\left\{|Q|^{-1/s'}\int_Q |F(x)|\,dx\right\}_{Q\in\mathbb D^0}\right\|_{\ell^{s,\infty}}
\leq C
\|F\|_{L^s(\mathbb R^N)}.
\]
\end{lemma}
\begin{proof}
Let $\mathcal E\subset\mathbb D^0$ be finite and set
\[
G_{\mathcal E}(x)
:=
\sum_{Q\in\mathcal E}|Q|^{-1/s'}\chi_Q(x).
\]
Since the dyadic cubes containing any fixed $x$ are totally ordered by inclusion, for each
$x\in\bigcup_{Q\in\mathcal E}Q$ we may choose a cube $Q_x\in\mathcal E$ of minimal volume such that $x\in Q_x$. Every $Q\in\mathcal E$ containing $x$ is then a dyadic ancestor of $Q_x$. Since the volume increases by the factor $2^N$ at each generation,
\[
G_{\mathcal E}(x)
=
\sum_{\substack{Q\in\mathcal E\\ x\in Q}}|Q|^{-1/s'}
\leq
|Q_x|^{-1/s'}
\sum_{k=0}^{\infty}2^{-kN/s'}
\lesssim_{N,s}
\sup_{Q\in\mathcal E}|Q|^{-1/s'}\chi_Q(x).
\]
Consequently,
\begin{align*}
\|G_{\mathcal E}\|_{L^{s'}(\mathbb R^N)}^{s'}
\lesssim_{N,s}
\int_{\mathbb R^N}
\sup_{Q\in\mathcal E}|Q|^{-1}\chi_Q(x)\,dx\leq
\int_{\mathbb R^N}
\sum_{Q\in\mathcal E}|Q|^{-1}\chi_Q(x)\,dx
=
\#\mathcal E.
\end{align*}
and hence,
\[
\|G_{\mathcal E}\|_{L^{s'}(\mathbb R^N)}
\lesssim_{N,s}
(\#\mathcal E)^{1/s'}.
\]
By H\"older's inequality,
\begin{align*}
\sum_{Q\in\mathcal E}|Q|^{-1/s'}\int_Q|F(x)|\,dx
=
\int_{\mathbb R^N}|F(x)|G_{\mathcal E}(x)\,dx\lesssim_{N,s}
\|F\|_{L^s(\mathbb R^N)}(\#\mathcal E)^{1/s'}.
\end{align*}

Fix $t>0$ and let $\mathcal E$ be any finite subcollection of the dyadic cubes satisfying
\[
|Q|^{-1/s'}\int_Q|F(x)|\,dx>t.
\]
Then
\[
t\,\#\mathcal E
<
\sum_{Q\in\mathcal E}|Q|^{-1/s'}\int_Q|F(x)|\,dx
\lesssim_{N,s}
\|F\|_{L^s(\mathbb R^N)}(\#\mathcal E)^{1/s'},
\]
so
\[
\#\mathcal E
\lesssim_{N,s}
t^{-s}\|F\|_{L^s(\mathbb R^N)}^s.
\]
Since $\mathcal E$ is arbitrary,
\[
\#\Bigg\{
Q\in\mathbb D^0:
|Q|^{-1/s'}\int_Q|F(x)|\,dx>t
\Bigg\}
\lesssim_{N,s}
t^{-s}\|F\|_{L^s(\mathbb R^N)}^s,
\qquad t>0.
\]
The asserted weak $\ell^s$ estimate follows. This completes the proof of Lemma~\ref{lem:Cwikel-dyadic-distribution}.
\end{proof}


\bibliographystyle{plain}
\bibliography{references}

\end{document}
